\documentclass{article}
\usepackage[margin=0.8in]{geometry}
\usepackage{amsmath,amssymb,amsthm,mathtools}
\usepackage{mathrsfs,bm,bbm}
\usepackage{graphicx,booktabs,array,multirow,tabularx,longtable}
\usepackage{enumitem}
\usepackage{xcolor}
\usepackage{algorithm,algpseudocode}
\usepackage{subcaption,tikz}
\usepackage{siunitx,cancel,accents}
\usepackage{microtype}
\usepackage[numbers,sort&compress]{natbib}
\usepackage[colorlinks=true,linkcolor=blue,citecolor=blue,urlcolor=blue]{hyperref}
\usepackage{cleveref}
\setlength{\emergencystretch}{2em}
\newtheorem{assumption}{Assumption}
\newtheorem{definition}{Definition}
\newtheorem{theorem}{Theorem}
\newtheorem{lemma}{Lemma}
\newtheorem{proposition}{Proposition}
\newtheorem{openendedquestion}{Open-ended Question}
\newtheorem{conjecture}{Conjecture}
\newcommand{\coreref}[1]{\ref{#1}}

\begin{document}

\sloppy

\section*{stat\_\allowbreak{}bdd\_\allowbreak{}uniform\_\allowbreak{}log\_\allowbreak{}penalty}

\noindent\textit{Specialization:} v1\par

\paragraph{TL;DR.} The paper has two co-main results.  First, on the full \(\mathcal P_{\mathrm{NP}}(L,q)\) law class for integer \(q\geq1\) and \(L\geq4\), every law-independent point-indexed distance family with Borel fixed sections has outer-expected uniform risk at least of order \((\log n/n)^{1/4}\); because this decision class strictly contains CTY's literal displayed one-common-map class, the same logarithmic converse follows there.  This first result is converse-only and makes no full-class achievability claim.  Second, on the distinct exact displayed \(L\)-uniformized Euclidean-distance, uniform-kernel class \(\mathcal P_{12}(p,\nu,L)\), an explicit winsorized, Gram-stabilized signed-distance local-polynomial estimator attains expected risk of order \(a_n=(\log n/n)^{1/4}\) for every \(L\geq4\), while a fixed-known-interior-boundary lower bound over arbitrary known-geometry point-indexed Borel sections yields the matched outer-expected minimax frontier for \(L\geq L_0(p)\).

\section{Project justification}

\paragraph{Gap.} Two distinct gaps are closed.  On the full \(\mathcal P_{\mathrm{NP}}(L,q)\) class, CTY Theorem~6 gives an \(n^{-1/4}\) lower rate for the literal common-map architecture and its discussion conjectures the logarithmic uniform penalty; no full-class upper conclusion follows from the additional geometry used by their upper theory.  Separately, on the displayed \(\mathcal P_{12}(p,\nu,L)\) specialization, CTY Theorem~1 supplies identification, Theorem~2 supplies a sequence-scoped first-order bias input, and Theorem~4(ii) supplies probability-order stochastic control, but CTY do not state an outer-expected lower/upper minimax equality on one common law, risk, and decision triple.

\paragraph{Niche.} The first construction attacks information loss from unsigned-distance compression while permitting an arbitrary law-independent Borel section at every boundary point.  The second uses known Euclidean assignment geometry and signed distances on the exact uniform-kernel \(\mathcal P_{12}\) specialization: its lower experiment fixes the support, assignment rectangle, and interior boundary, whereas its upper witness is an explicit winsorized, clipped, Gram-stabilized degree-\(p\) rule.

\paragraph{Fill.} For the first co-main result, a boundary angular hypercube and a coordinatewise overlap direct-product argument prove the \((\log n/n)^{1/4}\) converse throughout the full \(\mathcal P_{\mathrm{NP}}(L,q)\) class, first for arbitrary point-indexed Borel sections and then for the literal common-map class by strict inclusion.  For the second, a fixed-known-interior-boundary hypercube proves the same lower rate over all known-geometry point-indexed Borel sections, and direct comparison of a winsorized empirical score with the original population coefficient gives a uniform expected-risk upper bound.  Together they close the \(\mathcal P_{12}\) law-class, risk, and decision-class mismatch for \(L\geq L_0(p)\), while the upper witness remains valid for every \(L\geq4\).

\section{Related work}

Cattaneo, Titiunik, and Yu (2026), Theorem~6, is the direct comparator for the first co-main result: it studies the same \(\mathcal P_{\mathrm{NP}}(L,q)\) law class and literal displayed common-map distance architecture and gives an \(n^{-1/4}\) lower rate.  The present full-class converse proves the sharper \((\log n/n)^{1/4}\) lower rate over the strictly larger class of arbitrary law-independent point-indexed families with Borel fixed sections, and inclusion yields their common-map conclusion; no matching upper bound over all of \(\mathcal P_{\mathrm{NP}}\) is asserted.  The second co-main result has a different comparator and scope.  CTY Theorem~1 identifies the causal boundary curve, Theorem~2 gives the sequence-scoped first-order bias input, Theorem~4(ii) gives probability-order stochastic control under Assumptions~1--2, and Supplemental Appendix SA-8.3--SA-8.4 provides the maximal-inequality ingredients adapted here.  CTY Theorem~6 concerns the distinct support-boundary problem.  CTY do not state the present common-law-class outer-expected minimax equality on the exact displayed Euclidean-distance, uniform-kernel \(\mathcal P_{12}\) triple: the new fixed-known-interior-boundary lower bound covers arbitrary known-geometry point-indexed Borel sections, and the explicit winsorized, Gram-stabilized estimator supplies expected-risk achievability.

\section{Setup}

Target (estimand / structural parameter): First co-main target: \(\mu_P(\cdot)=\mathbb E_P[Y\mid X=\cdot]\big|_{\operatorname{bd}(\mathcal X_P)}\) on the exact \(\mathcal P_{\mathrm{NP}}(L,q)\) law class, for integer \(q\geq1\) and \(L\geq4\).  The result is a converse for arbitrary law-independent point-indexed Borel sections and, by inclusion, for the literal displayed CTY common-map class; no full-\(\mathcal P_{\mathrm{NP}}\) achievability claim is made.  Second co-main target: \(\tau_P(\cdot)=\mathbb E_P[Y(1)-Y(0)\mid X=\cdot]\big|_{\mathcal B_P}\) on the distinct exact displayed \(L\)-uniformized Euclidean-distance, uniform-kernel CTY Assumptions~1--2 plus Theorem-2-envelope class \(\mathcal P_{12}(p,\nu,L)\), with \(p\geq0\), \(q=p+1\geq1\), \(\nu\geq2\), and \(L\geq4\), identified by the signed-distance limits in CTY Theorem~1.  Its explicit expected-risk upper bound holds for every \(L\geq4\), while the matched outer-expected minimax conclusion additionally requires \(L\geq L_0(p)\)..
Primary functional / estimator / diagnostic: The note has two distinct, non-substitutable statistical functionals.  First, for each \(x\in\operatorname{bd}(\mathcal X_P)\), the observed input is \(V_{n,P}(x)=((Y_i,\lVert X_i-x\rVert_2):1\leq i\leq n)\).  The enlarged class \(\mathcal T_n^{\mathrm{PI}}\) contains all law-independent point-indexed families whose fixed-\(x\) sections are Borel.  The first co-main theorem proves \(R_n^{\mathrm{PI}}\gtrsim(\log n/n)^{1/4}\) on the full \(\mathcal P_{\mathrm{NP}}(L,q)\) class for integer \(q\geq1\) and \(L\geq4\).  This class strictly enlarges the literal common-map class displayed by CTY, \(\mathcal T=\{T_n(U_n(x)):T_n\in\mathcal T_n\}\), so inclusion gives the common-map converse on the same law class.  This contribution is converse-only and supplies no full-\(\mathcal P_{\mathrm{NP}}\) upper theorem.

Second, for \(P\in\mathcal P_{12}(p,\nu,L)\), where \(p\geq0\), \(q=p+1\geq1\), \(\nu\geq2\), and \(L\geq4\), the causal target is \(\tau_P(x)=\mathbb E_P[Y(1)-Y(0)\mid X=x]\), \(x\in\mathcal B_P\).  Under CTY Theorem~1 it equals the difference of the two one-sided signed-distance conditional-mean limits.  The estimator observes \(U^{\pm}_{n,P}(x)=((Y_i,D^{\pm}_{P,x,i}):1\leq i\leq n)\), may use the known geometry \(G_P\), and belongs to the law-independent sectionwise-Borel class \(\mathcal T_n^{12,\pm}\).  Its outer-expectation minimax functional is
\[
 R_n^{12,\pm}(p,\nu,L)=\inf_{\widehat\tau_n\in\mathcal T_n^{12,\pm}}
 \sup_{P\in\mathcal P_{12}(p,\nu,L)}
 \mathbb E_P^*\!\left[\sup_{x\in\mathcal B_P}|\widehat\tau_{n,P}(x)-\tau_P(x)|\right].
\]
The explicit winsorized, Gram-stabilized degree-\(p\) signed-distance local-polynomial estimator with \(h_n=a_n\), \(\psi_B(y)=\operatorname{sign}(y)\min\{|y|,B\}\), and \(B_n=a_n^{-1/3}\) attains the upper rate \(a_n=(\log n/n)^{1/4}\) for every \(p\geq0\), \(\nu\geq2\), and \(L\geq4\).  The fixed-known-interior-boundary converse over arbitrary known-geometry point-indexed Borel sections gives the matching lower rate when \(L\geq L_0(p)\).  Hence the exact matched outer-expected rate holds on this one common \(\mathcal P_{12}\) law, risk, and decision triple, not on all of \(\mathcal P_{\mathrm{NP}}\)..
\begin{itemize}
  \item \texttt{n} --- \texttt{sample size}; \(\mathbb{N}\); \texttt{index}
  \par\smallskip\noindent\emph{Definition.} \texttt{The number of independently sampled units.}
  \item \texttt{q} --- \texttt{regression smoothness order}; \(\mathbb{N}\); \texttt{regime parameter}
  \par\smallskip\noindent\emph{Definition.} The H\"older order of the single observed-outcome regression \(\mu_P\), with \(q\geq1\).
  \item \texttt{L} --- \texttt{CTY regularity envelope}; \([4,\infty)\); \texttt{law-\allowbreak{}class constant}
  \par\smallskip\noindent\emph{Definition.} A fixed envelope for compact support, score density, regression H\"older norm, and conditional variance in \(\mathcal P_{\mathrm{NP}}(L,q)\).
  \item \texttt{X} --- \texttt{bivariate random-\allowbreak{}design covariate}; \(\mathbb{R}^2\); \texttt{data component}
  \par\smallskip\noindent\emph{Definition.} \texttt{The observed location covariate whose support boundary indexes the regression target.}
  \item \texttt{Y} --- \texttt{observed outcome}; \(\mathbb{R}\); \texttt{data component}
  \par\smallskip\noindent\emph{Definition.} \texttt{The realized outcome.}
  \item \texttt{P} --- \texttt{data-\allowbreak{}generating law}; laws of \((Y,X)\); \texttt{model}
  \par\smallskip\noindent\emph{Definition.} \texttt{A member law for the observed outcome and bivariate random-\allowbreak{}design covariate.}
  \item \texttt{\textbackslash{}\allowbreak{}mathcal\{X\}\allowbreak{}\_\allowbreak{}P} --- \texttt{score support}; compact subsets of \(\mathbb{R}^2\); \texttt{geometric primitive}
  \par\smallskip\noindent\emph{Definition.} The support of \(X\) under \(P\).
  \item \texttt{f\_\allowbreak{}P} --- \texttt{score density}; \(\mathcal{X}_P\to(0,\infty)\); \texttt{nuisance}
  \par\smallskip\noindent\emph{Definition.} The Lebesgue density of \(X\) under \(P\).
  \item \texttt{a\_\allowbreak{}n} --- \texttt{first-\allowbreak{}order distance rate}; \texttt{positive sequences}; \texttt{frontier rate}
  \par\smallskip\noindent\emph{Definition.} \(a_n=(\log n/n)^{1/4}\).
  \item \texttt{\textbackslash{}\allowbreak{}mu\_\allowbreak{}P} --- \texttt{single observed-\allowbreak{}outcome regression}; \(\mathcal X_P\to\mathbb R\); \texttt{Cattaneo-\allowbreak{}-\allowbreak{}Titiunik-\allowbreak{}-\allowbreak{}Yu support-\allowbreak{}boundary target}
  \par\smallskip\noindent\emph{Definition.} \(\mu_P(x)=\mathbb E_P[Y\mid X=x]\).
  \item \texttt{\textbackslash{}\allowbreak{}sigma\_\allowbreak{}P\textasciicircum{}2} --- \texttt{single observed-\allowbreak{}outcome conditional variance}; \(\mathcal X_P\to(0,\infty)\); \texttt{CTY class nuisance}
  \par\smallskip\noindent\emph{Definition.} \(\sigma_P^2(x)=\operatorname{Var}_P(Y\mid X=x)\).
  \item \texttt{V\_\allowbreak{}\{n,\allowbreak{}P\}\allowbreak{}(x)} --- \texttt{unsigned-\allowbreak{}distance compressed sample}; \((\mathbb R^2)^n\); \texttt{CTY compressed observation}
  \par\smallskip\noindent\emph{Definition.} \(V_{n,P}(x)=((Y_i,\lVert X_i-x\rVert_2):1\leq i\leq n)\).
  \item \texttt{\textbackslash{}\allowbreak{}widehat\{\textbackslash{}\allowbreak{}mu\}\allowbreak{}\_\allowbreak{}n} --- \texttt{coordinatewise support-\allowbreak{}boundary regression estimator}; \(x\mapsto\widehat\mu_n(x)\); \texttt{rule in CTY's literal displayed common-\allowbreak{}map class}
  \par\smallskip\noindent\emph{Definition.} A generic rule in the literal common-map class displayed by CTY, \(\mathcal T=\{T_n(U_n(x)):T_n\in\mathcal T_n\}\), generated by one common law-independent measurable map \(T_n\) through \(\widehat\mu_n(x)=T_n(V_{n,P}(x))\) at every boundary point; \(x\) affects the rule only through the compressed sample.
  \item \texttt{\textbackslash{}\allowbreak{}mathcal\{T\}\allowbreak{}\textasciicircum{}\{\textbackslash{}\allowbreak{}mathrm\{NP\}\allowbreak{}\}\allowbreak{}\_\allowbreak{}n} --- \texttt{literal displayed CTY common-\allowbreak{}map coordinatewise distance decision class}; \texttt{families generated by one common measurable map}; \texttt{literal displayed CTY common-\allowbreak{}map comparator class}
  \par\smallskip\noindent\emph{Definition.} The internal class corresponding to the literal common-map class displayed by CTY, \(\mathcal T=\{T_n(U_n(x)):T_n\in\mathcal T_n\}\); its members use one common law-independent measurable map \(T_n\) at all boundary points, as constructed in \(\mathrm{def:cty\text{-}distance\text{-}decision\text{-}class}\).
  \item \texttt{\textbackslash{}\allowbreak{}mathcal\{P\}\allowbreak{}\_\allowbreak{}\{\textbackslash{}\allowbreak{}mathrm\{NP\}\allowbreak{}\}\allowbreak{}} --- \texttt{CTY nonparametric regression law class}; sets of laws \(P\); same-\(\mathcal P_{\mathrm{NP}}(L,q)\)-law-class comparator model
  \par\smallskip\noindent\emph{Definition.} The nonparametric law class constructed in \(\mathrm{def:cty\text{-}nonparametric\text{-}class}\) and used for the same-law-class comparison with CTY Theorem~6.
  \item \texttt{R\_\allowbreak{}n\textasciicircum{}\{\textbackslash{}\allowbreak{}mathrm\{NP\}\allowbreak{}\}\allowbreak{}} --- \texttt{literal-\allowbreak{}common-\allowbreak{}map coordinatewise distance minimax risk}; \([0,\infty)\); common-map frontier object on the same \(\mathcal P_{\mathrm{NP}}(L,q)\) law class
  \par\smallskip\noindent\emph{Definition.} The risk for the internal class corresponding to CTY's literal displayed common-map architecture, constructed in \(\mathrm{def:cty\text{-}distance\text{-}risk}\).
  \item \texttt{\textbackslash{}\allowbreak{}widehat\{\textbackslash{}\allowbreak{}mu\}\allowbreak{}\textasciicircum{}\{\textbackslash{}\allowbreak{}mathrm\{PI\}\allowbreak{}\}\allowbreak{}\_\allowbreak{}n} --- \texttt{point-\allowbreak{}indexed support-\allowbreak{}boundary regression estimator}; \(x\mapsto\widehat\mu_{n,P}^{\mathrm{PI}}(x)\); \texttt{sectionwise-\allowbreak{}Borel point-\allowbreak{}indexed decision rule}
  \par\smallskip\noindent\emph{Definition.} A generic point-indexed rule satisfying \(\widehat\mu_{n,P}^{\mathrm{PI}}(x)=T_{n,x}(V_{n,P}(x))\), where the law-independent family has a Borel measurable fixed-\(x\) section for every \(x\in\mathbb R^2\).
  \item \texttt{\textbackslash{}\allowbreak{}mathcal\{T\}\allowbreak{}\textasciicircum{}\{\textbackslash{}\allowbreak{}mathrm\{PI\}\allowbreak{}\}\allowbreak{}\_\allowbreak{}n} --- \texttt{point-\allowbreak{}indexed coordinatewise distance decision class}; \texttt{families with a Borel measurable section at each fixed point}; \texttt{strict enlargement of CTY's literal displayed common-\allowbreak{}map decision class}
  \par\smallskip\noindent\emph{Definition.} The sectionwise-Borel point-indexed decision class constructed in \(\mathrm{def:point\text{-}indexed\text{-}distance\text{-}decision\text{-}class}\); no joint measurability or other regularity in \(x\) is imposed.
  \item \texttt{R\_\allowbreak{}n\textasciicircum{}\{\textbackslash{}\allowbreak{}mathrm\{PI\}\allowbreak{}\}\allowbreak{}} --- \texttt{point-\allowbreak{}indexed distance minimax risk}; \texttt{extended nonnegative reals}; larger-decision-class frontier on the same \(\mathcal P_{\mathrm{NP}}(L,q)\) law class
  \par\smallskip\noindent\emph{Definition.} The outer-expectation minimax risk constructed in \(\mathrm{def:point\text{-}indexed\text{-}distance\text{-}risk}\).
  \item \texttt{p} --- \texttt{local-\allowbreak{}polynomial order}; \(\mathbb N_0\); \texttt{regular A1/\allowbreak{}A2 regime parameter}
  \par\smallskip\noindent\emph{Definition.} The signed-distance local-polynomial degree, with \(p\geq0\) and structural smoothness index \(q=p+1\geq1\).
  \item \texttt{\textbackslash{}\allowbreak{}nu} --- \texttt{conditional moment exponent offset}; \([2,\infty)\); \texttt{tail-\allowbreak{}regime parameter}
  \par\smallskip\noindent\emph{Definition.} The CTY raw conditional-moment exponent is \(2+\nu\).
  \item \texttt{h} --- \texttt{bandwidth}; \texttt{positive real sequences}; \texttt{estimator tuning parameter}
  \par\smallskip\noindent\emph{Definition.} \texttt{The signed-\allowbreak{}distance local-\allowbreak{}polynomial bandwidth.}
  \item \texttt{\textbackslash{}\allowbreak{}mathcal\{A\}\allowbreak{}\_\allowbreak{}\{t,\allowbreak{}P\}\allowbreak{}} --- \texttt{known assignment region}; Borel subsets of \(\mathcal X_P\); \texttt{assignment geometry}
  \par\smallskip\noindent\emph{Definition.} The known score region assigned to arm \(t\in\{0,1\}\).
  \item \texttt{T} --- \texttt{deterministic treatment indicator}; \(\{0,1\}\); \texttt{causal design component}
  \par\smallskip\noindent\emph{Definition.} \(T=\mathbf1\{X\in\mathcal A_{1,P}\}\).
  \item \texttt{Y(t)} --- potential outcome under arm \(t\); \(\mathbb R\); \texttt{causal outcome primitive}
  \par\smallskip\noindent\emph{Definition.} The potential outcome for \(t\in\{0,1\}\), with observed outcome \(Y=TY(1)+(1-T)Y(0)\).
  \item \texttt{\textbackslash{}\allowbreak{}mathcal\{B\}\allowbreak{}\_\allowbreak{}P} --- \texttt{known assignment boundary}; compact rectifiable subsets of \(\mathbb R^2\); \texttt{causal target index set}
  \par\smallskip\noindent\emph{Definition.} \(\mathcal B_P=\operatorname{bd}(\mathcal A_{0,P})\cap\operatorname{bd}(\mathcal A_{1,P})\subset\operatorname{int}(\mathcal X_P)\).
  \item \texttt{d\_\allowbreak{}P} --- \texttt{known design metric}; \(\mathbb R^2\times\mathbb R^2\to[0,\infty)\); \texttt{signed-\allowbreak{}distance geometry}
  \par\smallskip\noindent\emph{Definition.} For the exact class in this note, \(d_P(x,z)=\lVert x-z\rVert_2\).
  \item \texttt{K\_\allowbreak{}\textbackslash{}\allowbreak{}square} --- \texttt{uniform kernel}; \(\mathbb R\to\{0,1\}\); \texttt{fixed CTY kernel}
  \par\smallskip\noindent\emph{Definition.} \(K_\square(u)=\mathbf1\{|u|\leq1\}\).
  \item \texttt{G\_\allowbreak{}P} --- \texttt{known assignment geometry}; \texttt{design tuples}; \texttt{known estimator input}
  \par\smallskip\noindent\emph{Definition.} \(G_P=(\mathcal X_P,\mathcal A_{0,P},\mathcal A_{1,P},\mathcal B_P,d_P,K_\square)\).
  \item \texttt{\textbackslash{}\allowbreak{}mu\_\allowbreak{}\{t,\allowbreak{}P\}\allowbreak{}} --- \texttt{potential-\allowbreak{}outcome regression}; \(\mathcal X_P\to\mathbb R\); \texttt{A1 nuisance regression}
  \par\smallskip\noindent\emph{Definition.} \(\mu_{t,P}(x)=\mathbb E_P[Y(t)\mid X=x]\).
  \item \texttt{\textbackslash{}\allowbreak{}sigma\textasciicircum{}2\_\allowbreak{}\{t,\allowbreak{}P\}\allowbreak{}} --- \texttt{potential-\allowbreak{}outcome conditional variance}; \(\mathcal X_P\to(0,\infty)\); \texttt{A1 variance nuisance}
  \par\smallskip\noindent\emph{Definition.} \(\sigma^2_{t,P}(x)=\operatorname{Var}_P(Y(t)\mid X=x)\).
  \item \texttt{\textbackslash{}\allowbreak{}tau\_\allowbreak{}P} --- \texttt{boundary average treatment-\allowbreak{}effect curve}; \(\mathcal B_P\to\mathbb R\); \texttt{lead causal estimand}
  \par\smallskip\noindent\emph{Definition.} \(\tau_P(x)=\mathbb E_P[Y(1)-Y(0)\mid X=x]=\mu_{1,P}(x)-\mu_{0,P}(x)\).
  \item \texttt{r\_\allowbreak{}p} --- \texttt{polynomial basis}; \(\mathbb R\to\mathbb R^{p+1}\); \texttt{local-\allowbreak{}polynomial basis}
  \par\smallskip\noindent\emph{Definition.} \(r_p(u)=(1,u,\ldots,u^p)^\top\).
  \item \texttt{D\textasciicircum{}\{\textbackslash{}\allowbreak{}pm\}\allowbreak{}\_\allowbreak{}\{P,\allowbreak{}x\}\allowbreak{}} --- \texttt{known-\allowbreak{}geometry signed distance}; \(\mathbb R\); \texttt{compressed score}
  \par\smallskip\noindent\emph{Definition.} \(D^{\pm}_{P,x}=\{\mathbf1(X\in\mathcal A_{1,P})-\mathbf1(X\in\mathcal A_{0,P})\}d_P(X,x)\).
  \item \texttt{U\textasciicircum{}\{\textbackslash{}\allowbreak{}pm\}\allowbreak{}\_\allowbreak{}\{n,\allowbreak{}P\}\allowbreak{}(x)} --- \texttt{signed-\allowbreak{}distance compressed sample}; \((\mathbb R\times\mathbb R)^n\); \texttt{A1/\allowbreak{}A2 observed estimator input}
  \par\smallskip\noindent\emph{Definition.} \(U^{\pm}_{n,P}(x)=((Y_i,D^{\pm}_{P,x,i}):1\leq i\leq n)\).
  \item \texttt{\textbackslash{}\allowbreak{}Psi\_\allowbreak{}\{t,\allowbreak{}P,\allowbreak{}x\}\allowbreak{}(h)} --- \texttt{population signed-\allowbreak{}distance Gram matrix}; positive semidefinite \((p+1)\)-square matrices; \texttt{A2 stability functional}
  \par\smallskip\noindent\emph{Definition.} The population matrix displayed in \(\mathrm{def:cty\text{-}a1\text{-}a2\text{-}class}\).
  \item \texttt{\textbackslash{}\allowbreak{}mathcal\{P\}\allowbreak{}\_\allowbreak{}\{12\}\allowbreak{}} --- \texttt{uniformized CTY Assumptions 1-\allowbreak{}-\allowbreak{}2 plus Theorem-\allowbreak{}2-\allowbreak{}envelope law class}; sets of causal laws \(P\); \texttt{lead causal frontier model on the exact displayed CTY specialization}
  \par\smallskip\noindent\emph{Definition.} The exact displayed \(L\)-uniformized Euclidean-distance, uniform-kernel specialization constructed in \(\mathrm{def:cty\text{-}a1\text{-}a2\text{-}class}\), for \(p\geq0\), \(q=p+1\geq1\), and \(\nu\geq2\); it is not CTY's entire metric/kernel regime.
  \item \texttt{\textbackslash{}\allowbreak{}mathcal\{T\}\allowbreak{}\textasciicircum{}\{12,\allowbreak{}\textbackslash{}\allowbreak{}pm\}\allowbreak{}\_\allowbreak{}n} --- \texttt{known-\allowbreak{}geometry point-\allowbreak{}indexed signed-\allowbreak{}distance decision class}; \texttt{law-\allowbreak{}independent families with Borel fixed sections}; \texttt{lead minimax decision class}
  \par\smallskip\noindent\emph{Definition.} The sectionwise-Borel class constructed in \(\mathrm{def:cty\text{-}a1\text{-}a2\text{-}point\text{-}indexed\text{-}decision\text{-}class}\).
  \item \texttt{R\_\allowbreak{}n\textasciicircum{}\{12,\allowbreak{}\textbackslash{}\allowbreak{}pm\}\allowbreak{}} --- \texttt{A1/\allowbreak{}A2 signed-\allowbreak{}distance minimax risk}; \texttt{extended nonnegative reals}; \texttt{lead frontier object}
  \par\smallskip\noindent\emph{Definition.} The outer-expectation risk constructed in \(\mathrm{def:cty\text{-}a1\text{-}a2\text{-}outer\text{-}risk}\).
  \item \texttt{\textbackslash{}\allowbreak{}psi\_\allowbreak{}B} --- \texttt{winsorization map}; \(\mathbb R\to[-B,B]\); \texttt{heavy-\allowbreak{}tail stabilization}
  \par\smallskip\noindent\emph{Definition.} \(\psi_B(y)=\operatorname{sign}(y)\min\{|y|,B\}\) for \(B>0\).
  \item \texttt{\textbackslash{}\allowbreak{}widehat\{\textbackslash{}\allowbreak{}tau\}\allowbreak{}\textasciicircum{}\{12,\allowbreak{}\textbackslash{}\allowbreak{}mathrm\{LP\}\allowbreak{}\}\allowbreak{}\_\allowbreak{}\{n,\allowbreak{}h\}\allowbreak{}} --- \texttt{winsorized stabilized signed-\allowbreak{}distance local-\allowbreak{}polynomial estimator}; \(x\mapsto\widehat\tau^{12,\mathrm{LP}}_{n,h}(x)\); \texttt{upper-\allowbreak{}bound witness}
  \par\smallskip\noindent\emph{Definition.} The explicit winsorized, clipped, Gram-stabilized degree-\(p\) rule in \(\mathrm{def:cty\text{-}stabilized\text{-}local\text{-}polynomial\text{-}estimator}\), using winsorization level \(B(h)=h^{-1/3}\).
  \item \texttt{B\_\allowbreak{}n} --- \texttt{winsorization level}; \texttt{positive sequences}; \texttt{estimator tuning parameter}
  \par\smallskip\noindent\emph{Definition.} \(B_n=a_n^{-1/3}\).
\end{itemize}

\section{Assumptions}

\section{Definitions}

\begin{definition}[def:frontier-rates]\label{def:frontier-rates}
\par\smallskip\noindent\emph{Defined object.} \texttt{a\_\allowbreak{}n}
\par\smallskip\noindent\emph{Construction.} \[a_n=(\log n/n)^{1/4}.\]
\end{definition}

\begin{definition}[def:cty-nonparametric-class]\label{def:cty-nonparametric-class}
\par\smallskip\noindent\emph{Defined object.} \texttt{\textbackslash{}\allowbreak{}mathcal\{P\}\allowbreak{}\_\allowbreak{}\{\textbackslash{}\allowbreak{}mathrm\{NP\}\allowbreak{}\}\allowbreak{}}
\par\smallskip\noindent\emph{Construction.} \[\mathcal P_{\mathrm{NP}}(L,q)=\bigl\{P:\text{conditions (i)--(iv) below hold}\bigr\}.\]
Condition \(\mathrm{(i)}\): for every \(n\), \((Y_i,X_i)_{i=1}^n\) are independently and identically distributed under \(P\) in \(\mathbb R\times\mathbb R^2\).  Condition \(\mathrm{(ii)}\): \(X\) has a Lebesgue density \(f_P\), continuous on its compact support \(\mathcal X_P\subseteq[-L,L]^2\), with \(L^{-1}\leq\inf_{x\in\mathcal X_P}f_P(x)\leq\sup_{x\in\mathcal X_P}f_P(x)\leq L\), and \(\operatorname{bd}(\mathcal X_P)\) is a rectifiable curve.  Condition \(\mathrm{(iii)}\): \(\mu_P(x)=\mathbb E_P[Y\mid X=x]\) belongs to \(\mathcal H^q(\mathcal X_P,L)\).  Condition \(\mathrm{(iv)}\): \(\sigma_P^2(x)=\operatorname{Var}_P(Y\mid X=x)\) is continuous on \(\mathcal X_P\) and satisfies \(L^{-1}\leq\sigma_P^2(x)\leq L\) for every \(x\in\mathcal X_P\).
\end{definition}

\begin{definition}[def:cty-distance-data]\label{def:cty-distance-data}
\par\smallskip\noindent\emph{Defined object.} \texttt{V\_\allowbreak{}\{n,\allowbreak{}P\}\allowbreak{}(x)}
\par\smallskip\noindent\emph{Construction.} \[V_{n,P}(x)=((Y_i,\lVert X_i-x\rVert_2):1\leq i\leq n),\qquad x\in\operatorname{bd}(\mathcal X_P).\]
\end{definition}

\begin{definition}[def:cty-distance-decision-class]\label{def:cty-distance-decision-class}
\par\smallskip\noindent\emph{Defined object.} \texttt{\textbackslash{}\allowbreak{}mathcal\{T\}\allowbreak{}\textasciicircum{}\{\textbackslash{}\allowbreak{}mathrm\{NP\}\allowbreak{}\}\allowbreak{}\_\allowbreak{}n}
\par\smallskip\noindent\emph{Construction.} \[
 \mathcal T_n^{\mathrm{NP}}
 =\left\{\widehat\mu_n:\begin{array}{l}
 \text{there is one measurable map }T_n:(\mathbb R\times[0,\infty))^n\to\mathbb R
 \text{ selected independently of }P,\\
 \widehat\mu_n(x)=T_n(V_{n,P}(x))
 \text{ for every }x\in\operatorname{bd}(\mathcal X_P)
 \end{array}\right\}.
\]
\par\smallskip\noindent\emph{Inputs.} \texttt{def:\allowbreak{}cty-\allowbreak{}distance-\allowbreak{}data}.
\end{definition}

\begin{definition}[def:cty-distance-risk]\label{def:cty-distance-risk}
\par\smallskip\noindent\emph{Defined object.} \texttt{R\_\allowbreak{}n\textasciicircum{}\{\textbackslash{}\allowbreak{}mathrm\{NP\}\allowbreak{}\}\allowbreak{}}
\par\smallskip\noindent\emph{Construction.} \[R_n^{\mathrm{NP}}=\inf_{\widehat\mu_n\in\mathcal T_n^{\mathrm{NP}}}\sup_{P\in\mathcal P_{\mathrm{NP}}(L,q)}\mathbb E_P\!\left[\sup_{x\in\operatorname{bd}(\mathcal X_P)}|\widehat\mu_n(x)-\mu_P(x)|\right].\]
\par\smallskip\noindent\emph{Inputs.} \texttt{def:\allowbreak{}cty-\allowbreak{}nonparametric-\allowbreak{}class}; \texttt{def:\allowbreak{}cty-\allowbreak{}distance-\allowbreak{}decision-\allowbreak{}class}.
\end{definition}

\begin{definition}[def:point-indexed-distance-decision-class]\label{def:point-indexed-distance-decision-class}
\par\smallskip\noindent\emph{Defined object.} \texttt{\textbackslash{}\allowbreak{}mathcal\{T\}\allowbreak{}\textasciicircum{}\{\textbackslash{}\allowbreak{}mathrm\{PI\}\allowbreak{}\}\allowbreak{}\_\allowbreak{}n}
\par\smallskip\noindent\emph{Construction.} \[
 \mathcal T_n^{\mathrm{PI}}
 =\left\{\widehat\mu_n^{\mathrm{PI}}:\begin{array}{l}
 \text{there is a family }\{T_{n,x}:x\in\mathbb R^2\}\text{ selected independently of }P,\\
 T_{n,x}:(\mathbb R\times[0,\infty))^n\to\mathbb R
 \text{ is Borel measurable for every fixed }x,\\
 \widehat\mu_{n,P}^{\mathrm{PI}}(x)=T_{n,x}(V_{n,P}(x))
 \text{ for every }P\in\mathcal P_{\mathrm{NP}}(L,q)\\
 \text{and every }x\in\operatorname{bd}(\mathcal X_P)
 \end{array}\right\}.
\]
No joint measurability, continuity, common modulus, or other regularity of \(x\mapsto T_{n,x}\) is required.
\par\smallskip\noindent\emph{Inputs.} \texttt{def:\allowbreak{}cty-\allowbreak{}distance-\allowbreak{}data}; \texttt{def:\allowbreak{}cty-\allowbreak{}nonparametric-\allowbreak{}class}.
\end{definition}

\begin{definition}[def:point-indexed-distance-risk]\label{def:point-indexed-distance-risk}
\par\smallskip\noindent\emph{Defined object.} \texttt{R\_\allowbreak{}n\textasciicircum{}\{\textbackslash{}\allowbreak{}mathrm\{PI\}\allowbreak{}\}\allowbreak{}}
\par\smallskip\noindent\emph{Construction.} For an extended nonnegative function \(Z\) on the sample space, define its outer expectation by
\[
 \mathbb E_P^*[Z]
 =\inf\left\{\mathbb E_P[W]:
 W\text{ is an extended nonnegative measurable function and }W\geq Z
 \text{ pointwise}\right\}.
\]
Then
\[
 R_n^{\mathrm{PI}}
 =\inf_{\widehat\mu_n^{\mathrm{PI}}\in\mathcal T_n^{\mathrm{PI}}}
 \sup_{P\in\mathcal P_{\mathrm{NP}}(L,q)}
 \mathbb E_P^*\!\left[
 \sup_{x\in\operatorname{bd}(\mathcal X_P)}
 \left|\widehat\mu_{n,P}^{\mathrm{PI}}(x)-\mu_P(x)\right|
 \right].
\]
\par\smallskip\noindent\emph{Inputs.} \texttt{def:\allowbreak{}cty-\allowbreak{}nonparametric-\allowbreak{}class}; \texttt{def:\allowbreak{}point-\allowbreak{}indexed-\allowbreak{}distance-\allowbreak{}decision-\allowbreak{}class}.
\end{definition}

\begin{definition}[def:cty-uniform-kernel]\label{def:cty-uniform-kernel}
\par\smallskip\noindent\emph{Defined object.} \texttt{K\_\allowbreak{}\textbackslash{}\allowbreak{}square}
\par\smallskip\noindent\emph{Construction.} Define the fixed uniform kernel by
\[
 K_\square(u)=\mathbf 1\{|u|\leq1\},\qquad u\in\mathbb R.
\]
\end{definition}

\begin{definition}[def:cty-known-geometry]\label{def:cty-known-geometry}
\par\smallskip\noindent\emph{Defined object.} \texttt{G\_\allowbreak{}P}
\par\smallskip\noindent\emph{Construction.} For a member law \(P\), the known design object is
\[
 G_P=(\mathcal X_P,\mathcal A_{0,P},\mathcal A_{1,P},\mathcal B_P,d_P,K_\square),
\]
where \(\mathcal A_{0,P}\) and \(\mathcal A_{1,P}\) are a Borel partition of \(\mathcal X_P\), \(\mathcal B_P=\operatorname{bd}(\mathcal A_{0,P})\cap\operatorname{bd}(\mathcal A_{1,P})\subset\operatorname{int}(\mathcal X_P)\), and the fixed metric is \(d_P(x,z)=\lVert x-z\rVert_2\).
\par\smallskip\noindent\emph{Inputs.} \texttt{def:\allowbreak{}cty-\allowbreak{}uniform-\allowbreak{}kernel}.
\end{definition}

\begin{definition}[def:cty-known-geometry-signed-distance-data]\label{def:cty-known-geometry-signed-distance-data}
\par\smallskip\noindent\emph{Defined object.} \texttt{U\textasciicircum{}\{\textbackslash{}\allowbreak{}pm\}\allowbreak{}\_\allowbreak{}\{n,\allowbreak{}P\}\allowbreak{}(x)}
\par\smallskip\noindent\emph{Construction.} For \(x\in\mathcal B_P\), define
\[
 D^{\pm}_{P,x}=\{\mathbf 1(X\in\mathcal A_{1,P})-\mathbf 1(X\in\mathcal A_{0,P})\}\lVert X-x\rVert_2,
 \qquad
 U^{\pm}_{n,P}(x)=((Y_i,D^{\pm}_{P,x,i}):1\leq i\leq n).
\]
\par\smallskip\noindent\emph{Inputs.} \texttt{def:\allowbreak{}cty-\allowbreak{}known-\allowbreak{}geometry}.
\end{definition}

\begin{definition}[def:cty-a1-a2-class]\label{def:cty-a1-a2-class}
\par\smallskip\noindent\emph{Defined object.} \texttt{\textbackslash{}\allowbreak{}mathcal\{P\}\allowbreak{}\_\allowbreak{}\{12\}\allowbreak{}}
\par\smallskip\noindent\emph{Construction.} For integers \(p\geq0\), real numbers \(\nu\geq2\), and \(L\geq4\), let \(\mathcal P_{12}(p,\nu,L)\) be the set of laws \(P\) satisfying all the following member properties.

\(\mathrm{(i)}\) For every \(n\), the vectors \((Y_i(0),Y_i(1),X_i)\), \(1\leq i\leq n\), are independently and identically distributed; \(T=\mathbf 1\{X\in\mathcal A_{1,P}\}\); and \(Y=TY(1)+(1-T)Y(0)\).

\(\mathrm{(ii)}\) The support has the rectangular form \(\mathcal X_P=[\underline x_P,\overline x_P]^2\subseteq[-L,L]^2\).  The variable \(X\) has a Lebesgue density \(f_P\), continuous on \(\mathcal X_P\), such that \(L^{-1}\leq f_P(x)\leq L\) for all \(x\in\mathcal X_P\).

\(\mathrm{(iii)}\) For each \(t\in\{0,1\}\), \(\mu_{t,P}(x)=\mathbb E_P[Y(t)\mid X=x]\) has a \(C^{p+1}\) extension to an open neighborhood of \(\mathcal X_P\), and
\[
 \max_{0\leq|\alpha|\leq p}\sup_{x\in\mathcal X_P}|\partial^\alpha\mu_{t,P}(x)|
 +\max_{0\leq|\alpha|\leq p}\sup_{x\neq z\in\mathcal X_P}
 { |\partial^\alpha\mu_{t,P}(x)-\partial^\alpha\mu_{t,P}(z)|\over\lVert x-z\rVert_2}
 \leq L.
\]

\(\mathrm{(iv)}\) For each \(t\), \(\sigma^2_{t,P}(x)=\operatorname{Var}_P(Y(t)\mid X=x)\) is continuous and satisfies \(L^{-1}\leq\sigma^2_{t,P}(x)\leq L\) on \(\mathcal X_P\).

\(\mathrm{(v)}\) For each \(t\) and every \(x\in\mathcal X_P\),
\[
 \mathbb E_P[|Y(t)|^{2+\nu}\mid X=x]\leq L.
\]

\(\mathrm{(vi)}\) The sets \(\mathcal A_{0,P}\) and \(\mathcal A_{1,P}\) form a Borel partition of \(\mathcal X_P\), and their known common interface \(\mathcal B_P\subset\operatorname{int}(\mathcal X_P)\) is a compact rectifiable curve satisfying \(L^{-1}\leq\mathcal H^1(\mathcal B_P)\leq L\).

\(\mathrm{(vii)}\) The known metric is exactly Euclidean, \(d_P(x,z)=\lVert x-z\rVert_2\), and the known kernel is \(K_\square\).  The collection
\[
 \{\{z\in\mathcal X_P:\lVert z-x\rVert_2\leq r\}:x\in\mathcal X_P, r\geq0\}
\]
is a VC class with VC index at most four.

\(\mathrm{(viii)}\) Put \(r_p(u)=(1,u,\ldots,u^p)^\top\), \(I_1=[0,\infty)\), \(I_0=(-\infty,0)\), and
\[
 \Psi_{t,P,x}(h)=\mathbb E_P\!\left[
 {1\over h^2}\mathbf 1\{D^{\pm}_{P,x}\in I_t\}K_\square(D^{\pm}_{P,x}/h)
 r_p(D^{\pm}_{P,x}/h)r_p(D^{\pm}_{P,x}/h)^\top
 \right].
\]
For every \(0<h\leq L^{-1}\), every \(x\in\mathcal B_P\), and both \(t\),
\[
 \lambda_{\min}(\Psi_{t,P,x}(h))\geq L^{-1}.
\]

\(\mathrm{(ix)}\) For every \(0<h\leq L^{-1}\), every \(x\in\mathcal B_P\), and both \(t\),
\[
 {1\over h^2}\int_{\mathcal A_{t,P}}
 K_\square(\{(2t-1)\lVert z-x\rVert_2\}/h)\,dz\geq L^{-1}.
\]

\(\mathrm{(x)}\) For every \(x\in\mathcal B_P\), both \(t\), and every \(0<s\leq L^{-1}\), the one-dimensional Hausdorff integral of \(f_P\) over
\[
 \{z\in\mathcal A_{t,P}:\lVert z-x\rVert_2=s\}
\]
is finite and strictly positive.

This is the exact displayed \(L\)-uniformized Euclidean-distance, uniform-kernel specialization of CTY Assumptions~1--2, augmented by the displayed Theorem-2 envelope restrictions and quantitative small-bandwidth thresholds.  It is not CTY's entire metric/kernel regime.
\par\smallskip\noindent\emph{Inputs.} \texttt{def:\allowbreak{}cty-\allowbreak{}uniform-\allowbreak{}kernel}; \texttt{def:\allowbreak{}cty-\allowbreak{}known-\allowbreak{}geometry}; \texttt{def:\allowbreak{}cty-\allowbreak{}known-\allowbreak{}geometry-\allowbreak{}signed-\allowbreak{}distance-\allowbreak{}data}.
\end{definition}

\begin{definition}[def:cty-a1-a2-point-indexed-decision-class]\label{def:cty-a1-a2-point-indexed-decision-class}
\par\smallskip\noindent\emph{Defined object.} \texttt{\textbackslash{}\allowbreak{}mathcal\{T\}\allowbreak{}\textasciicircum{}\{12,\allowbreak{}\textbackslash{}\allowbreak{}pm\}\allowbreak{}\_\allowbreak{}n}
\par\smallskip\noindent\emph{Construction.} Let \(\mathcal T^{12,\pm}_n\) be the class of rules for which there is one collection
\[
 \{T_{n,G,x}:G\text{ is an admissible known geometry},\ x\in\mathbb R^2\}
\]
selected independently of the unknown member law, such that every fixed \((G,x)\) section is a Borel map from \((\mathbb R\times\mathbb R)^n\) to \(\mathbb R\) and
\[
 \widehat\tau_{n,P}(x)=T_{n,G_P,x}(U^{\pm}_{n,P}(x)),
 \qquad P\in\mathcal P_{12}(p,\nu,L),\quad p\geq0,\quad x\in\mathcal B_P.
\]
No joint measurability, continuity, common modulus, or other regularity in \(x\) is imposed.
\par\smallskip\noindent\emph{Inputs.} \texttt{def:\allowbreak{}cty-\allowbreak{}known-\allowbreak{}geometry}; \texttt{def:\allowbreak{}cty-\allowbreak{}known-\allowbreak{}geometry-\allowbreak{}signed-\allowbreak{}distance-\allowbreak{}data}; \texttt{def:\allowbreak{}cty-\allowbreak{}a1-\allowbreak{}a2-\allowbreak{}class}.
\end{definition}

\begin{definition}[def:cty-a1-a2-outer-risk]\label{def:cty-a1-a2-outer-risk}
\par\smallskip\noindent\emph{Defined object.} \texttt{R\_\allowbreak{}n\textasciicircum{}\{12,\allowbreak{}\textbackslash{}\allowbreak{}pm\}\allowbreak{}}
\par\smallskip\noindent\emph{Construction.} For an extended nonnegative function \(Z\), put
\[
 \mathbb E_P^*[Z]=\inf\{\mathbb E_P[W]:W\text{ is extended nonnegative and measurable, and }W\geq Z\text{ pointwise}\}.
\]
For \(p\geq0\), define
\[
 R^{12,\pm}_n(p,\nu,L)=
 \inf_{\widehat\tau_n\in\mathcal T^{12,\pm}_n}
 \sup_{P\in\mathcal P_{12}(p,\nu,L)}
 \mathbb E_P^*\!\left[
 \sup_{x\in\mathcal B_P}|\widehat\tau_{n,P}(x)-\tau_P(x)|
 \right].
\]
\par\smallskip\noindent\emph{Inputs.} \texttt{def:\allowbreak{}cty-\allowbreak{}a1-\allowbreak{}a2-\allowbreak{}class}; \texttt{def:\allowbreak{}cty-\allowbreak{}a1-\allowbreak{}a2-\allowbreak{}point-\allowbreak{}indexed-\allowbreak{}decision-\allowbreak{}class}.
\end{definition}

\begin{definition}[def:cty-stabilized-local-polynomial-estimator]\label{def:cty-stabilized-local-polynomial-estimator}
\par\smallskip\noindent\emph{Defined object.} \texttt{\textbackslash{}\allowbreak{}widehat\{\textbackslash{}\allowbreak{}tau\}\allowbreak{}\textasciicircum{}\{12,\allowbreak{}\textbackslash{}\allowbreak{}mathrm\{LP\}\allowbreak{}\}\allowbreak{}\_\allowbreak{}\{n,\allowbreak{}h\}\allowbreak{}}
\par\smallskip\noindent\emph{Construction.} For \(B>0\), define \(\psi_B(y)=\operatorname{sign}(y)\min\{|y|,B\}\).  For each \(t\in\{0,1\}\), \(x\in\mathcal B_P\), and \(h>0\), let
\[
 \widehat\Psi_{t,x}(h)={1\over nh^2}\sum_{i=1}^n
 \mathbf 1\{D^{\pm}_{P,x,i}\in I_t\}K_\square(D^{\pm}_{P,x,i}/h)
 r_p(D^{\pm}_{P,x,i}/h)r_p(D^{\pm}_{P,x,i}/h)^\top
\]
and, with \(B(h)=h^{-1/3}\), let
\[
 \widehat m^B_{t,x}(h)={1\over nh^2}\sum_{i=1}^n
 \mathbf 1\{D^{\pm}_{P,x,i}\in I_t\}K_\square(D^{\pm}_{P,x,i}/h)
 r_p(D^{\pm}_{P,x,i}/h)\psi_{B(h)}(Y_i).
\]
If \(\lambda_{\min}(\widehat\Psi_{t,x}(h))\geq(2L)^{-1}\), set \(\widehat\beta^B_{t,x}(h)=\widehat\Psi_{t,x}(h)^{-1}\widehat m^B_{t,x}(h)\); otherwise set \(\widehat\beta^B_{t,x}(h)=0\).  With \(e_1=(1,0,\ldots,0)^\top\), define
\[
 \widehat\tau^{12,\mathrm{LP}}_{n,h}(x)=
 \Pi_{[-2L,2L]}\!\left\{e_1^\top(\widehat\beta^B_{1,x}(h)-\widehat\beta^B_{0,x}(h))\right\}.
\]
At the theorem's bandwidth \(h_n=a_n\), the winsorization level is \(B_n=B(h_n)=a_n^{-1/3}\).  This is a fixed, law-independent, jointly Borel member of \(\mathcal T^{12,\pm}_n\) for every \(p\geq0\).
\par\smallskip\noindent\emph{Inputs.} \texttt{def:\allowbreak{}frontier-\allowbreak{}rates}; \texttt{def:\allowbreak{}cty-\allowbreak{}uniform-\allowbreak{}kernel}; \texttt{def:\allowbreak{}cty-\allowbreak{}known-\allowbreak{}geometry-\allowbreak{}signed-\allowbreak{}distance-\allowbreak{}data}; \texttt{def:\allowbreak{}cty-\allowbreak{}a1-\allowbreak{}a2-\allowbreak{}class}; \texttt{def:\allowbreak{}cty-\allowbreak{}a1-\allowbreak{}a2-\allowbreak{}point-\allowbreak{}indexed-\allowbreak{}decision-\allowbreak{}class}.
\end{definition}

\section{Main results}

\begin{lemma}[lem:coordinatewise-overlap-direct-product]\label{lem:coordinatewise-overlap-direct-product}
Let all spaces carrying \(Z_0,Z_1,\ldots,Z_M,S_1,\ldots,S_M\) be standard Borel.  Let \(\Omega_1,\ldots,\Omega_M\) be independent uniform binary variables.  Conditionally on them, let \(Z_0,Z_1,\ldots,Z_M\) be independent, with the law of \(Z_j\) depending only on \(\Omega_j\) for \(1\leq j\leq M\), and with the law of \(Z_0\) independent of \((\Omega_1,\ldots,\Omega_M)\).  Suppose decoder \(j\) is measurable with respect to \((S_j,Z_{-j},Z_0)\), where \(S_j\) is a measurable compression of \(Z_j\), and write \(Q_{j,b}=\mathcal L(S_j\mid\Omega_j=b)\) and \(\rho_j=1-\lVert Q_{j,0}-Q_{j,1}\rVert_{\mathrm{TV}}\).  Then every such decentralized decoder obeys
\[
 \Pr\{\widehat\Omega_j=\Omega_j\text{ for every }j\}
 \leq\frac12\left\{1+\prod_{j=1}^M(1-\rho_j)\right\}.
\]
In particular, if \(\operatorname{KL}(Q_{j,0},Q_{j,1})\leq\kappa\log M\) for all \(j\) and a fixed \(\kappa<1\), then
\[
 \Pr\{\text{at least one decoding error}\}
 \geq\frac12\left[1-\exp\{-M^{1-\kappa}/2\}\right].
\]
Along any sequence with \(M=M_n\to\infty\) and fixed \(\kappa<1\), this lower bound is \(1/2-o(1)\).
\end{lemma}
\begin{proof}
For each \(j\), dominate \(Q_{j,0}\) and \(Q_{j,1}\) by \(Q_{j,0}+Q_{j,1}\), and write the resulting densities as \(q_{j,0}\) and \(q_{j,1}\).  The common subprobability measure has density \(q_{j,0}\wedge q_{j,1}\) and mass
\[
 \rho_j=\int(q_{j,0}\wedge q_{j,1})
 =1-\lVert Q_{j,0}-Q_{j,1}\rVert_{\mathrm{TV}}.
\]
Construct an equivalent version of the experiment as follows.  Independently across coordinates, draw an ambiguity indicator \(A_j\), independent of \(\Omega_j\), with \(\Pr(A_j=1)=\rho_j\).  On \(\{A_j=1\}\), draw \(S_j\) from the normalized common submeasure, independently of \(\Omega_j\).  On \(\{A_j=0\}\), draw \(S_j\) from the normalized residual measure belonging to \(\Omega_j\).  When a normalizing mass is zero, assign an arbitrary probability kernel to the corresponding probability-zero branch.  The resulting conditional law of \(S_j\) given \(\Omega_j=b\) is exactly \(Q_{j,b}\).

Because the spaces are standard Borel and \(S_j\) is a measurable compression of \(Z_j\), a regular conditional distribution of \(Z_j\) given \((S_j,\Omega_j)\) exists.  Sampling from those kernels reconstructs the prescribed conditional law of every \(Z_j\), and sampling \(Z_0\) independently from its bit-independent law reconstructs the joint experiment by the assumed conditional product structure.  Standard-Borel kernel randomization permits all these draws to be realized using latent uniforms independent of the bits.

Condition on the ambiguity indicators, all latent randomization other than the bits, and all bits outside
\[
 \mathcal A=\{j:A_j=1\}.
\]
Suppose \(\mathcal A\neq\varnothing\).  The bit configurations on \(\mathcal A\) for which all decoders are correct form an independent set in the \(|\mathcal A|\)-dimensional binary hypercube.  Indeed, take two configurations differing only in coordinate \(j\in\mathcal A\).  The decoder for coordinate \(j\) receives the same \(S_j\), because the common-submeasure branch is independent of \(\Omega_j\); it also receives the same \(Z_{-j}\) and \(Z_0\), because every other bit and all latent randomization are fixed.  Its input is therefore identical although the correct value of \(\Omega_j\) is reversed, so the decoder cannot be correct at both endpoints of that hypercube edge.  Pairing vertices along any fixed coordinate shows that an independent set has at most \(2^{|\mathcal A|-1}\) vertices.  The ambiguous bits remain independent and uniform under the conditioning.  Hence the conditional probability of simultaneous success is at most \(1/2\) on \(\{\mathcal A\neq\varnothing\}\), while it is at most one on \(\{\mathcal A=\varnothing\}\).  Averaging and using independence of the \(A_j\)'s gives
\[
\begin{aligned}
 \Pr\{\widehat\Omega_j=\Omega_j\text{ for every }j\}
 &\leq \Pr(\mathcal A=\varnothing)+\frac12\Pr(\mathcal A\neq\varnothing)\\
 &=\frac12\left\{1+\prod_{j=1}^M(1-\rho_j)\right\}.
\end{aligned}
\]

It remains to derive the finite-\(M\) relative-entropy consequence.  Let
\[
 H_j=\int\sqrt{q_{j,0}q_{j,1}}
\]
be the Hellinger affinity.  Jensen's inequality under \(Q_{j,0}\) yields
\[
 \operatorname{KL}(Q_{j,0},Q_{j,1})\geq-2\log H_j.
\]
The elementary affinity bound for total variation follows from Cauchy--Schwarz:
\[
\begin{aligned}
 \lVert Q_{j,0}-Q_{j,1}\rVert_{\mathrm{TV}}
 &=\frac12\int|\sqrt{q_{j,0}}-\sqrt{q_{j,1}}|
             (\sqrt{q_{j,0}}+\sqrt{q_{j,1}})\\
 &\leq\sqrt{1-H_j^2}.
\end{aligned}
\]
Consequently,
\[
 \rho_j
 \geq1-\sqrt{1-e^{-\operatorname{KL}(Q_{j,0},Q_{j,1})}}
 \geq\frac12e^{-\operatorname{KL}(Q_{j,0},Q_{j,1})},
\]
where the last inequality uses \(1-\sqrt{1-u}=u/(1+\sqrt{1-u})\geq u/2\) for \(0\leq u\leq1\).  Under the asserted entropy bound,
\[
 \rho_j\geq\frac12M^{-\kappa}.
\]
Therefore \(1-v\leq e^{-v}\) gives
\[
 \prod_{j=1}^M(1-\rho_j)
 \leq\exp\!\left\{-\sum_{j=1}^M\rho_j\right\}
 \leq\exp\{-M^{1-\kappa}/2\}.
\]
Subtracting the simultaneous-success bound from one proves
\[
 \Pr\{\text{at least one decoding error}\}
 \geq\frac12\left[1-\exp\{-M^{1-\kappa}/2\}\right].
\]
If \(M=M_n\to\infty\) and \(\kappa<1\) is fixed, then \(M_n^{1-\kappa}\to\infty\), so the exponential term is \(o(1)\), proving the final assertion.
\end{proof}
\par\smallskip\noindent \textit{Justification.} The finite-\(M\) overlap decomposition gives the exact simultaneous-error certificate \(\frac12[1-\exp\{-M^{1-\kappa}/2\}]\); its asymptotic \(1/2-o(1)\) consequence is invoked only after the packing verifies \(M=M_n\to\infty\). \textit{Gap.} The lemma allows a different measurable decoder at every coordinate and therefore supplies the multiple-coordinate step needed for arbitrary point-indexed sections. \textit{Consumer.} The point-indexed converse applies the lemma with \(M=M_n\) and \(\kappa=2\alpha<1\).

\begin{lemma}[lem:cty-support-boundary-angular-packing]\label{lem:cty-support-boundary-angular-packing}
For every \(q\geq1\) and \(L\geq4\), there are constants \(c_0,c_1>0\) and \(0<\alpha<1/8\), depending only on \(q\) and \(L\), such that, for all sufficiently large \(n\), there exist \(M_n\geq c_0a_n^{-1/q}\) points \(x_1,\ldots,x_{M_n}\) on the boundary of the square \(S=[-1/2,1/2]^2\), separated by at least a constant multiple of \(a_n^{1/q}\), a radius \(w_n\asymp a_n^{1/q}\), and laws \(P_\omega\in\mathcal P_{\mathrm{NP}}(L,q)\), indexed by \(\omega\in\{0,1\}^{M_n}\), with the following properties.  Every law has support \(S\), a continuous density \(f_{P_\omega}\) satisfying \(L^{-1}\leq f_{P_\omega}\leq L\), a regression \(\mu_{P_\omega}\) in the stated \(q\)-H\"older ball, and a continuous conditional variance \(\sigma^2_{P_\omega}\) satisfying \(L^{-1}\leq\sigma^2_{P_\omega}\leq L\).  The boundary \(\operatorname{bd}(S)\) is rectifiable and the observations are independently and identically distributed.  The half-discs
\[
 C_j=B_2(x_j,w_n)\cap S,\qquad 1\leq j\leq M_n,
\]
are pairwise disjoint.  The subprobability law \(P_\omega\{(Y,X)\in\cdot,\ X\in C_j\}\) depends on \(\omega\) only through \(\omega_j\); the subprobability law \(P_\omega\{(Y,X)\in\cdot,\ X\in S\setminus\bigcup_{k=1}^{M_n}C_k\}\) is independent of \(\omega\); and \(P_\omega(X\in C_j)=m_n\) for one common number \(m_n\) independent of \(j\) and \(\omega\).  There are numbers \(u_{j,0},u_{j,1}\) such that
\[
 \mu_{P_\omega}(x_j)=u_{j,\omega_j},
 \qquad |u_{j,1}-u_{j,0}|\geq c_1a_n.
\]
If \(\omega^{(j)}\) differs from \(\omega\) only in coordinate \(j\), then the laws of \(\lVert X-x_j\rVert_2\) agree, the conditional radial outcome laws agree outside radii of order \(a_n\), and
\[
 \operatorname{KL}\!\left(\mathcal L_{P_\omega}\{V_{n,P_\omega}(x_j)\},
 \mathcal L_{P_{\omega^{(j)}}}\{V_{n,P_{\omega^{(j)}}}(x_j)\}\right)
 \leq\alpha\log M_n.
\]
\end{lemma}
\begin{proof}
Put \(S=[-1/2,1/2]^2\), and place all centers on a fixed middle subinterval of the lower edge of \(S\).  If \(x_j\) is such a center, then, for every sufficiently small \(r>0\), the part of the circle of radius \(r\) lying in \(S\) is
\[
 x_j+r e_\theta,
 \qquad e_\theta=(\cos\theta,\sin\theta),
 \qquad 0\leq\theta\leq\pi.
\]
We use
\[
 \int_0^\pi\cos\theta\,d\theta=0,
 \qquad
 \int_0^\pi\cos^2\theta\,d\theta=\frac\pi2.
\]

Choose \(\psi\in C_c^\infty((-1,1))\) with \(0\leq\psi\leq1\) and \(\psi(0)=1\), and set
\[
 \phi(z)=\psi(\lVert z\rVert_2^2),
 \qquad
 \varphi(u)=\psi(u^2).
\]
For constants \(d,d_1>0\), to be chosen below, define
\[
 \Delta=d a_n,
 \qquad
 w=(d_1\Delta)^{1/q}.
\]
Choose a maximal grid of centers on the indicated middle subinterval with separation \(3w\).  For all sufficiently large \(n\), the half-discs \(C_j=B_2(x_j,w)\cap S\) are pairwise disjoint and avoid the corners.  Elementary interval packing gives
\[
 M_n\asymp w^{-1},
 \qquad
 M_n\geq c_0a_n^{-1/q}.
\]
Thus \(w_n=w\asymp a_n^{1/q}\), and the separation is a constant multiple of \(a_n^{1/q}\).  Write \(r_j(z)=\lVert z-x_j\rVert_2\) and \(\phi_j(z)=\phi((z-x_j)/w)\).

Fix a sufficiently small slope \(b>0\).  For \(\omega\in\{0,1\}^{M_n}\), define
\[
 p_\omega(z)=\frac12+bz_1+\Delta\sum_{j=1}^{M_n}\omega_j\phi_j(z).
\]
For every integer \(0\leq k\leq q\), scaling gives
\[
 \lVert D^k(\Delta\phi_j)\rVert_\infty
 \leq C_{\phi,q}\Delta w^{-k}.
\]
For \(k=q\), the right-hand side is \(C_{\phi,q}/d_1\); for \(k<q\), it converges to zero.  The supports are disjoint, so summation does not enlarge these suprema.  Choose \(b\) small, next \(d_1\) large, and finally \(d\) small.  The affine part and scaled bumps then put every \(p_\omega\) in \(\mathcal H^q(S,L)\), and, for all sufficiently large \(n\),
\[
 \frac14\leq p_\omega(z)\leq\frac34.
\]

Let \(\chi\in C^\infty(\mathbb R)\) satisfy \(0\leq\chi\leq1\), \(\chi(v)=0\) for \(v\leq1\), and \(\chi(v)=1\) for \(v\geq2\).  Fix \(c_A\geq8\), increasing \(d_1\) if necessary so that \(2c_A\Delta/b<w\) eventually, including when \(q=1\), and define
\[
 A_n(r)=-\frac{2\Delta\varphi(r/w)}{br}
 \chi\!\left(\frac{br}{c_A\Delta}\right),
 \qquad A_n(0)=0.
\]
The cutoff makes \(A_n\) identically zero in a neighborhood of zero, while compact support of \(\varphi(r/w)\) makes it vanish smoothly before \(r=w\).  On the region where it can be nonzero, \(br>c_A\Delta\), and hence
\[
 |A_n(r)|\leq\frac{2}{c_A}\leq\frac14.
\]
Define on \(S\)
\[
 f_\omega(z)=1+
 \sum_{j=1}^{M_n}\omega_j A_n(r_j(z))
 \frac{z_1-(x_j)_1}{r_j(z)}\mathbf 1\{r_j(z)<w\},
\]
with each summand defined to be zero at its center.  Disjointness makes the formula unambiguous.  The two cutoff properties show that \(f_\omega\) is continuous on \(S\), and
\[
 \frac34\leq f_\omega(z)\leq\frac54.
\]
For each affected half-disc, polar integration gives
\[
 \int_0^w A_n(r)r\,dr\int_0^\pi\cos\theta\,d\theta=0.
\]
Therefore \(\int_S f_\omega(z)\,dz=1\).  Positivity makes \(S\) the support.  Since \(L\geq4\), the displayed density interval is contained in \([L^{-1},L]\), and \(S\subset[-L,L]^2\) is compact with rectifiable boundary.

Conditional on \(X=z\), let
\[
 Y=B+\varepsilon,
 \qquad
 B\sim\operatorname{Bernoulli}(p_\omega(z)),
 \qquad
 \varepsilon\sim N(0,1),
\]
where \(B\) and \(\varepsilon\) are conditionally independent.  Then
\[
 \mu_{P_\omega}(z)=p_\omega(z),
 \qquad
 \sigma_{P_\omega}^2(z)=1+p_\omega(z)\{1-p_\omega(z)\}.
\]
The variance is continuous and belongs to \([1,5/4]\subseteq[L^{-1},L]\).  Draw the \(n\) observations independently from this common law.  Thus every clause of \(\mathrm{def:cty\text{-}nonparametric\text{-}class}\) holds, so every \(P_\omega\) belongs to \(\mathcal P_{\mathrm{NP}}(L,q)\).

Because the half-discs are disjoint, on \(C_j\) both \(p_\omega\) and \(f_\omega\), and hence the joint law of \((Y,X)\) restricted to \(X\in C_j\), depend on \(\omega\) only through \(\omega_j\).  Outside \(\bigcup_k C_k\), all perturbations vanish, so the restricted joint law is independent of \(\omega\).  The zero angular integral and congruence of the half-discs give
\[
 P_\omega(X\in C_j)=\int_{C_j}f_\omega(z)\,dz
 =|C_j|=\frac{\pi w^2}{2}=:m_n.
\]
Disjointness and \(\phi_j(x_j)=1\) yield
\[
 \mu_{P_\omega}(x_j)=\frac12+b(x_j)_1+\Delta\omega_j
 =:u_{j,\omega_j},
 \qquad |u_{j,1}-u_{j,0}|=\Delta=d a_n.
\]

It remains to establish the compressed-experiment properties.  Compare adjacent vertices, orienting the notation first so that \(\omega_j=1\) and \(\omega_j^{(j)}=0\).  At every \(0<r<w\), the angular density correction has integral zero, so the two radial densities agree and equal \(\pi r\).  For \(r\geq w\), the adjacent score densities are pointwise identical.  Thus the complete laws of \(\lVert X-x_j\rVert_2\) agree.

For \(0<r<w\), the difference of the outcome-weighted angular integrals is
\[
\begin{aligned}
 &\int_0^\pi p_\omega(x_j+re_\theta)
 \{1+A_n(r)\cos\theta\}\,d\theta
 -\int_0^\pi p_{\omega^{(j)}}(x_j+re_\theta)\,d\theta\\
 &\qquad=\pi\Delta\varphi(r/w)+\frac\pi2brA_n(r).
\end{aligned}
\]
If \(br\geq2c_A\Delta\), then \(\chi=1\) and the last display is zero.  For \(r\geq w\), both adjacent perturbations vanish.  Conditional on a radius, an angular mixture of Bernoulli variables is Bernoulli with the angularly averaged success probability, and convolution with the common Gaussian preserves equality.  Therefore versions of the conditional radial outcome laws agree outside
\[
 0\leq r<\frac{2c_A\Delta}{b},
\]
an interval whose length is of order \(a_n\).

On this exceptional interval the common radial density is bounded by \(Cr\), both averaged Bernoulli probabilities lie in \([1/4,3/4]\), and their difference is at most \(C\Delta\).  The inequality \(\log t\leq t-1\), applied to the two Bernoulli atoms, gives
\[
 \operatorname{kl}(p,p')
 \leq\frac{(p-p')^2}{p'(1-p')}
 \leq C\Delta^2.
\]
The log-sum inequality shows that convolution with the common Gaussian cannot increase relative entropy.  Disintegration with respect to the common radial marginal consequently gives, for one compressed observation,
\[
 \operatorname{KL}(P^{\mathrm{rad}}_{\omega,j},P^{\mathrm{rad}}_{\omega^{(j)},j})
 \leq C\Delta^2\int_0^{2c_A\Delta/b}r\,dr
 \leq C'\Delta^4.
\]
The argument is unchanged after reversing the adjacent orientation.  Tensorization for the independent product gives
\[
 \operatorname{KL}\!\left(
 \mathcal L_{P_\omega}\{V_{n,P_\omega}(x_j)\},
 \mathcal L_{P_{\omega^{(j)}}}\{V_{n,P_{\omega^{(j)}}}(x_j)\}
 \right)
 \leq C'n\Delta^4=C'd^4\log n.
\]
Because \(M_n\asymp w^{-1}\),
\[
 \log M_n=\left(\frac{1}{4q}+o(1)\right)\log n.
\]
Choose \(d>0\) small enough that the divergence is at most \(\alpha\log M_n\) for one fixed \(0<\alpha<1/8\), without disturbing the earlier envelope choices.  This proves every asserted property.
\end{proof}
\par\smallskip\noindent \textit{Justification.} The construction is the same-law-class witness: it places separated boundary alternatives in \(\mathcal P_{\mathrm{NP}}(L,q)\) and proves cellwise locality, value separation, common cell mass, and the compressed Kullback--Leibler bound. \textit{Gap.} Its role is confined to the explicit square-support finite submodel; the outer-expectation minimax quantifiers and sectionwise estimator class are handled in the converse theorem. \textit{Consumer.} It supplies the finite hypercube and local divergence bound used by the point-indexed theorem and, through inclusion, by the CTY corollary.

\begin{theorem}[thm:cty-same-class-log-converse]\label{thm:cty-same-class-log-converse}
For every \(q\geq1\) and \(L\geq4\), there is a constant \(c_{\mathrm{CTY}}=c_{\mathrm{CTY}}(q,L)>0\) such that
\[
 \liminf_{n\to\infty}
 \left(\frac{n}{\log n}\right)^{1/4}
 \inf_{\widehat\mu_n\in\mathcal T_n^{\mathrm{NP}}}
 \sup_{P\in\mathcal P_{\mathrm{NP}}(L,q)}
 \mathbb E_P\!\left[
 \sup_{x\in\operatorname{bd}(\mathcal X_P)}
 |\widehat\mu_n(x)-\mu_P(x)|
 \right]
 =\liminf_{n\to\infty}\frac{R_n^{\mathrm{NP}}}{a_n}
 \geq c_{\mathrm{CTY}}.
\]
Thus the logarithmic penalty is necessary in the coordinatewise distance decision class and on the same nonparametric law class used by Cattaneo, Titiunik, and Yu, Theorem~6; the result is not a subclass transfer from the two-sided-thickness treatment-effect model.
\end{theorem}
\begin{proof}
Fix \(q\geq1\) and \(L\geq4\).  For a fixed \(n\geq2\), take any \(\widehat\mu_n\in\mathcal T_n^{\mathrm{NP}}\).  By \(\mathrm{def:cty\text{-}distance\text{-}decision\text{-}class}\), there is one law-independent Borel map
\[
 T_n:(\mathbb R\times[0,\infty))^n\longrightarrow\mathbb R
\]
such that \(\widehat\mu_n(x)=T_n(V_{n,P}(x))\) at every boundary point.  Equivalently, define \(T_{n,x}=T_n\) for every \(x\in\mathbb R^2\).  By \(\mathrm{lem:common\text{-}map\text{-}strict\text{-}in\text{-}point\text{-}indexed}\),
\[
 \mathcal T_n^{\mathrm{NP}}\subseteq\mathcal T_n^{\mathrm{PI}}.
\]

We next verify that the ordinary completed expectation in the CTY risk agrees with outer expectation for a common-map rule.  Fix \(P\in\mathcal P_{\mathrm{NP}}(L,q)\), write a sample point as \(\mathbf z=((y_i,\xi_i):1\leq i\leq n)\), and put
\[
 D_n(x,\mathbf z)=((y_i,\lVert\xi_i-x\rVert_2):1\leq i\leq n),
 \qquad
 \ell_{P,n}(x,\mathbf z)
 =\left|T_n(D_n(x,\mathbf z))-\mu_P(x)\right|.
\]
The map \((x,\mathbf z)\mapsto D_n(x,\mathbf z)\) is Borel, \(T_n\) is Borel, and \(\mu_P\) is continuous because it lies in \(\mathcal H^q(\mathcal X_P,L)\).  Hence \(\ell_{P,n}\) is jointly Borel on \(\operatorname{bd}(\mathcal X_P)\times(\mathbb R\times\mathbb R^2)^n\).  The boundary is compact because \(\mathcal X_P\) is compact.  Consequently, for every real \(t\),
\[
 \left\{\mathbf z:
   \sup_{x\in\operatorname{bd}(\mathcal X_P)}
   \ell_{P,n}(x,\mathbf z)>t
 \right\}
 =\operatorname{proj}_{\mathbf z}
 \left\{(x,\mathbf z):
   x\in\operatorname{bd}(\mathcal X_P),\ 
   \ell_{P,n}(x,\mathbf z)>t
 \right\}.
\]
The set before projection is Borel in a product of standard Borel spaces, so its projection is analytic and belongs to the completion of every probability measure on the sample space.  Thus the boundary supremum is upper semianalytic and universally measurable.  Its completed expectation equals its outer expectation:
\[
 \mathbb E_P\!\left[
   \sup_{x\in\operatorname{bd}(\mathcal X_P)}
   |\widehat\mu_n(x)-\mu_P(x)|
 \right]
 =
 \mathbb E_P^*\!\left[
   \sup_{x\in\operatorname{bd}(\mathcal X_P)}
   |\widehat\mu_n(x)-\mu_P(x)|
 \right].
\]
This argument uses joint measurability only for the common-map rule.

The two minimax risks use the same law class.  The last equality and inclusion of decision classes give
\[
\begin{aligned}
 R_n^{\mathrm{NP}}
 &=\inf_{\widehat\mu_n\in\mathcal T_n^{\mathrm{NP}}}
   \sup_{P\in\mathcal P_{\mathrm{NP}}(L,q)}
   \mathbb E_P^*\!\left[
     \sup_{x\in\operatorname{bd}(\mathcal X_P)}
     |\widehat\mu_n(x)-\mu_P(x)|
   \right]\\
 &\geq R_n^{\mathrm{PI}}.
\end{aligned}
\]
The theorem \(\mathrm{thm:point\text{-}indexed\text{-}distance\text{-}log\text{-}converse}\) supplies \(c_{\mathrm{PI}}(q,L)>0\) such that
\[
 \liminf_{n\to\infty}\frac{R_n^{\mathrm{PI}}}{a_n}
 \geq c_{\mathrm{PI}}(q,L).
\]
Monotonicity of \(\liminf\) gives the same lower bound for \(R_n^{\mathrm{NP}}/a_n\).  Finally, \(a_n=(\log n/n)^{1/4}\) by \(\mathrm{def:frontier\text{-}rates}\), and \(R_n^{\mathrm{NP}}\) is the displayed infimum--supremum by \(\mathrm{def:cty\text{-}distance\text{-}risk}\).  Taking \(c_{\mathrm{CTY}}=c_{\mathrm{PI}}\) proves the claim on the exact same class.
\end{proof}
\par\smallskip\noindent \textit{Justification.} The result for the literal common-map class displayed by CTY follows from \(R_n^{\mathrm{NP}}\geq R_n^{\mathrm{PI}}\) and the point-indexed theorem.  Here ``same-class'' refers to the same \(\mathcal P_{\mathrm{NP}}(L,q)\) law class.  For common Borel maps the boundary-loss supremum is universally measurable, so ordinary completed expectation equals outer expectation. \textit{Gap.} It remains distinct from the new causal A1/A2 family. \textit{Consumer.} This resolves the printed CTY logarithmic lower-bound conjecture for the literal displayed common-map class in the stated regime \(q\in\mathbb N\), \(q\geq1\), and \(L\geq4\).

\begin{lemma}[lem:common-map-strict-in-point-indexed]\label{lem:common-map-strict-in-point-indexed}
For every \(n\geq1\), \(q\geq1\), and \(L\geq4\),
\[
 \mathcal T_n^{\mathrm{NP}}\subsetneq\mathcal T_n^{\mathrm{PI}}.
\]
\end{lemma}
\begin{proof}
Take a rule in \(\mathcal T_n^{\mathrm{NP}}\), generated by one law-independent Borel map
\[
 T_n:(\mathbb R\times[0,\infty))^n\longrightarrow\mathbb R.
\]
Set \(T_{n,x}=T_n\) for every \(x\in\mathbb R^2\).  Each fixed section is Borel measurable, the family is selected independently of \(P\), and \(T_{n,x}(V_{n,P}(x))=T_n(V_{n,P}(x))\) at every boundary point.  Thus \(\mathcal T_n^{\mathrm{NP}}\subseteq\mathcal T_n^{\mathrm{PI}}\).

For strictness, consider the law-independent family
\[
 T_{n,x}(v)=x_1,
 \qquad x=(x_1,x_2)\in\mathbb R^2.
\]
For every fixed \(x\), this is a constant, hence Borel measurable, function of \(v\), so it defines a member of \(\mathcal T_n^{\mathrm{PI}}\).

Let \(P_0\) be the law under which \(X\) is uniform on \(S=[-1/2,1/2]^2\) and \(Y\sim N(0,1)\) is independent of \(X\).  Its density is one, its support boundary is rectifiable, its regression is zero, and its conditional variance is one.  Hence \(P_0\in\mathcal P_{\mathrm{NP}}(L,q)\) for every \(q\geq1\) and \(L\geq4\).  At
\[
 x_-=(-1/2,0),
 \qquad
 x_+=(1/2,0),
\]
reflection of the square in the vertical axis gives
\[
 \mathcal L_{P_0}\{V_{n,P_0}(x_-)\}
 =\mathcal L_{P_0}\{V_{n,P_0}(x_+)\}.
\]
If this point-indexed family were generated by one common Borel map, applying it to these identically distributed inputs would produce outputs with the same law.  The required outputs are instead the distinct constants
\[
 T_{n,x_-}(V_{n,P_0}(x_-))=-\frac12,
 \qquad
 T_{n,x_+}(V_{n,P_0}(x_+))=\frac12.
\]
This contradiction proves strict inclusion.
\end{proof}
\par\smallskip\noindent \textit{Justification.} The reflection witness proves that arbitrary Borel fixed sections genuinely enlarge the literal common-map class displayed by CTY, rather than merely reparameterizing it. \textit{Gap.} The witness requires no joint measurability: the family \(T_{n,x}(v)=x_1\) has Borel fixed sections, while equal reflected compressed laws rule out representation by one common map. \textit{Consumer.} It certifies the strict decision-class strengthening and the inclusion used for the corollary concerning CTY's literal displayed common-map class.

\begin{theorem}[thm:point-indexed-distance-log-converse]\label{thm:point-indexed-distance-log-converse}
For every \(q\geq1\) and \(L\geq4\), there is a constant \(c_{\mathrm{PI}}=c_{\mathrm{PI}}(q,L)>0\) such that
\[
 \liminf_{n\to\infty}
 \left(\frac{n}{\log n}\right)^{1/4}
 \inf_{\widehat\mu_n^{\mathrm{PI}}\in\mathcal T_n^{\mathrm{PI}}}
 \sup_{P\in\mathcal P_{\mathrm{NP}}(L,q)}
 \mathbb E_P^*\!\left[
 \sup_{x\in\operatorname{bd}(\mathcal X_P)}
 \left|\widehat\mu_{n,P}^{\mathrm{PI}}(x)-\mu_P(x)\right|
 \right]
 =\liminf_{n\to\infty}\frac{R_n^{\mathrm{PI}}}{a_n}
 \geq c_{\mathrm{PI}}.
\]
Here \(\mathcal T_n^{\mathrm{PI}}\) permits an arbitrary law-independent family \(\{T_{n,x}:x\in\mathbb R^2\}\) whose fixed-\(x\) sections are Borel measurable; no joint measurability, continuity, or common modulus in \(x\) is imposed, and \(\mathbb E_P^*\) is outer expectation.  Thus the logarithmic penalty remains necessary in a decision class strictly larger than the literal common-map class displayed by CTY, \(\mathcal T=\{T_n(U_n(x)):T_n\in\mathcal T_n\}\).  The theorem is a lower bound on the same nonparametric law class \(\mathcal P_{\mathrm{NP}}(L,q)\), not an Assumption~2 upper-bound result or a matched two-sided frontier.
\end{theorem}
\begin{proof}
Fix \(q\geq1\) and \(L\geq4\).  For every sufficiently large \(n\), take the family supplied by \(\mathrm{lem:cty\text{-}support\text{-}boundary\text{-}angular\text{-}packing}\).  Write \(M=M_n\), put \(\Delta=c_1a_n\), and let \(C_j=B_2(x_j,w_n)\cap S\).  The packing lemma gives
\[
 \mu_{P_\omega}(x_j)=u_{j,\omega_j},
 \qquad |u_{j,1}-u_{j,0}|\geq\Delta,
 \qquad M_n\geq c_0a_n^{-1/q}\longrightarrow\infty.
\]

Let \(\widehat\mu_n^{\mathrm{PI}}\in\mathcal T_n^{\mathrm{PI}}\) be arbitrary, generated by law-independent Borel fixed sections \(T_{n,x}\).  No joint measurability in \(x\) will be used.  Under \(P_\omega\), generate a Poisson random measure on \(\mathbb R\times S\times(0,1)\) with intensity \(2nP_\omega(dy,dx)du\), where \(u\) is an independent uniform mark.  Its total count \(N\) has law \(\operatorname{Pois}(2n)\).  If \(N\geq n\), retain the \(n\) observations with smallest marks, order them by their marks, and apply each fixed section to its compressed sample; if \(N<n\), return zero.  Conditional on \(N\geq n\), the retained ordered observations have law \(P_\omega^{\otimes n}\).

Put independent uniform bits \(\Omega_1,\ldots,\Omega_M\) on the vertices.  Let \(Z_j\) be the complete marked restriction of the Poisson random measure to \(C_j\), and let \(Z_0\) be its restriction to the complement of the cells.  Spaces of finite counting measures over standard Borel spaces are standard Borel.  Conditional on the bits, the restrictions are independent; cellwise locality makes the law of \(Z_j\) depend only on \(\Omega_j\), and the law of \(Z_0\) is bit-independent.

For each \(j\), let \(S_j\) retain from \(Z_j\) only every point's mark, outcome, and distance from \(x_j\).  The value \(\widehat\mu_n^{\mathrm{Pois}}(x_j)\) is measurable with respect to \((S_j,Z_{-j},Z_0)\): the off-cell raw blocks reveal every off-cell distance from \(x_j\), \(S_j\) supplies the in-cell compressed observations, the marks recover their order, and \(T_{n,x_j}\) is Borel.  This verification is coordinatewise and requires no regularity of \(x\mapsto T_{n,x}\).

Let \(\nu_{j,b}\) be the one-observation subprobability law of the outcome--distance pair restricted to \(C_j\).  It has common mass \(m_n\).  The complete adjacent compressed laws have disjoint-support form \(\nu_{j,b}+\lambda_j\), with \(\lambda_j\) common and supported on radii at least \(w_n\).  Define
\[
 D_j=\int\log\!\left(\frac{d\nu_{j,0}}{d\nu_{j,1}}\right)d\nu_{j,0}.
\]
The packing divergence bound and tensorization imply
\[
 nD_j\leq\alpha\log M.
\]
For \(Q_{j,b}=\mathcal L(S_j\mid\Omega_j=b)\), the cell count is \(K_j\sim\operatorname{Pois}(2nm_n)\) under either bit.  Conditional on \(K_j=k\), the compressed cell points are independent with law \(\nu_{j,b}/m_n\), and the marks contribute no entropy.  Therefore
\[
 \operatorname{KL}(Q_{j,0},Q_{j,1})
 =\mathbb E[K_j]\frac{D_j}{m_n}
 =2nD_j
 \leq2\alpha\log M.
\]
The reverse orientation is identical.  Since \(\alpha<1/8\), \(\kappa=2\alpha<1\).

Decode \(\Omega_j\) by thresholding \(\widehat\mu_n^{\mathrm{Pois}}(x_j)\) at the midpoint of \(u_{j,0}\) and \(u_{j,1}\), reversing the inequality if necessary.  By \(\mathrm{lem:coordinatewise\text{-}overlap\text{-}direct\text{-}product}\),
\[
 \Pr\{\text{at least one coordinate is decoded incorrectly}\}
 \geq\frac12\left[1-\exp\{-M_n^{1-2\alpha}/2\}\right].
\]
Define the measurable finite-coordinate Poissonized loss
\[
 L_{n,\omega}^{\mathrm{fin,Pois}}
 =\max_{1\leq j\leq M_n}
 \left|\widehat\mu_n^{\mathrm{Pois}}(x_j)-\mu_{P_\omega}(x_j)\right|.
\]
Any decoder error forces this loss to be at least \(\Delta/2\).  Averaging over the bit prior gives
\[
 \sup_\omega\mathbb E_{P_\omega}^{\mathrm{Pois}}
 [L_{n,\omega}^{\mathrm{fin,Pois}}]
 \geq\frac\Delta4
 \left[1-\exp\{-M_n^{1-2\alpha}/2\}\right].
\]

Now define the analogous loss \(L_{n,\omega}^{\mathrm{fin}}\) for the original fixed-size rule.  Since \(1/4\leq\mu_{P_\omega}\leq3/4\), the zero default has loss at most one on \(\{N<n\}\), while on \(\{N\geq n\}\) the retained observations have the original product law.  Hence
\[
 \mathbb E_{P_\omega}^{\mathrm{Pois}}[L_{n,\omega}^{\mathrm{fin,Pois}}]
 \leq \mathbb E_{P_\omega}[L_{n,\omega}^{\mathrm{fin}}]
 +\Pr\{\operatorname{Pois}(2n)<n\}.
\]
For \(N\sim\operatorname{Pois}(2n)\), exponential Markov with \(t=\log2\) gives
\[
 \Pr(N<n)\leq\exp\{-n(1-\log2)\}=o(a_n).
\]
Since \(M_n\to\infty\) and \(\Delta=c_1a_n\), it follows that
\[
 \sup_\omega\mathbb E_{P_\omega}[L_{n,\omega}^{\mathrm{fin}}]\geq ca_n
\]
for all sufficiently large \(n\).

Pointwise,
\[
 L_{n,\omega}^{\mathrm{fin}}
 \leq\sup_{x\in\operatorname{bd}(S)}
 \left|\widehat\mu_{n,P_\omega}^{\mathrm{PI}}(x)-\mu_{P_\omega}(x)\right|.
\]
Every measurable majorant of the possibly nonmeasurable supremum also majorizes the measurable finite maximum.  The definition of outer expectation therefore yields the same inequality after taking the corresponding expectations.  Since every packing vertex lies in \(\mathcal P_{\mathrm{NP}}(L,q)\) and the rule was arbitrary,
\[
 R_n^{\mathrm{PI}}\geq ca_n
\]
for all sufficiently large \(n\).  Divide by \(a_n=(\log n/n)^{1/4}\) and take lower limits.  Strict enlargement of the common-map class follows from \(\mathrm{lem:common\text{-}map\text{-}strict\text{-}in\text{-}point\text{-}indexed}\).
\end{proof}
\par\smallskip\noindent \textit{Justification.} This is the strengthened nonparametric headline: on the same \(\mathcal P_{\mathrm{NP}}(L,q)\) law class, it allows every fixed-point estimator section to be an arbitrary law-independent Borel map. \textit{Gap.} No full-class upper theorem is inferred from the separate causal result. \textit{Consumer.} It supplies the larger-decision-class minimax converse and implies the converse for the literal common-map class displayed by CTY, \(\mathcal T=\{T_n(U_n(x)):T_n\in\mathcal T_n\}\), by inclusion.

\begin{lemma}[lem:cty-distance-identification]\label{lem:cty-distance-identification}
Under CTY Assumptions~1(i)--(iii) and 2(i), (ii), and (v), the causal boundary treatment-effect curve is identified by the signed-distance conditional limits:
\[
 \tau_P(x)=\lim_{r\downarrow0}\mathbb E_P[Y\mid D^{\pm}_{P,x}=r]
 -\lim_{r\uparrow0}\mathbb E_P[Y\mid D^{\pm}_{P,x}=r],
 \qquad x\in\mathcal B_P.
\]
\end{lemma}
\par\smallskip\noindent \textit{Justification.} Identification is given by the design and cited separately from the new estimation frontier. \textit{Gap.} The causal target must not be collapsed by definition into the distance-based regression functional. \textit{Consumer.} The expected upper proposition and paper-level estimand metadata use this identifying equality.

\begin{lemma}[lem:cty-uniform-first-order-bias]\label{lem:cty-uniform-first-order-bias}
Fix an integer \(p\geq0\), a number \(\nu\geq2\), and \(L\geq4\).  For \(P\in\mathcal P_{12}(p,\nu,L)\), \(t\in\{0,1\}\), \(x\in\mathcal B_P\), and \(h>0\), let \(\beta^*_{t,P,x}(h)\in\mathbb R^{p+1}\) solve the original, unwinsorized population normal equations
\[
 \Psi_{t,P,x}(h)\beta^*_{t,P,x}(h)
 =\mathbb E_P\!\left[
 {1\over h^2}\mathbf1\{D^{\pm}_{P,x}\in I_t\}K_\square(D^{\pm}_{P,x}/h)
 r_p(D^{\pm}_{P,x}/h)Y
 \right],
\]
and put \(e_1=(1,0,\ldots,0)^\top\).  There is a constant \(C=C(p,L)<\infty\) such that, for every deterministic sequence \((h_n)_{n\geq1}\) satisfying \(h_n>0\), \(h_n\downarrow0\), and \(nh_n^2\to\infty\), for all sufficiently large \(n\),
\[
 \sup_{P\in\mathcal P_{12}(p,\nu,L)}
 \sup_{x\in\mathcal B_P}
 \left|
 e_1^\top\{\beta^*_{1,P,x}(h_n)-\beta^*_{0,P,x}(h_n)\}-\tau_P(x)
 \right|
 \leq Ch_n.
\]
\end{lemma}
\par\smallskip\noindent \textit{Justification.} This is the published first-order contrast-bias input, retained as a cited leaf in its source-supported sequential regime. \textit{Gap.} CTY Theorem~2 states the uniform conclusion along deterministic bandwidth sequences satisfying \(h_n\downarrow0\) and \(nh_n^2\to\infty\); the upper proposition verifies both conditions at \(h_n=a_n\). \textit{Consumer.} The explicit estimator upper proposition combines this sequential contrast-bias bound with the expected Gram and winsorized-score bounds.

\begin{lemma}[lem:cty-expected-local-polynomial-maximal-bounds]\label{lem:cty-expected-local-polynomial-maximal-bounds}
For the degree-\(p\) signed-distance local-polynomial score and Gram processes with the fixed uniform kernel, the expected uniform deviations over \(P\in\mathcal P_{12}(p,\nu,L)\), both arms, and \(x\in\mathcal B_P\) satisfy
\[
 \mathbb E_P^*\!\left[\sup_{t,x}\lVert\widehat\Psi_{t,x}(h)-\Psi_{t,P,x}(h)\rVert\right]
 \leq C\sqrt{{\log(1/h)\over nh^2}},
\]
and
\[
 \mathbb E_P^*\!\left[\sup_{t,x}\lVert\widehat O_{t,x}(h)\rVert\right]
 \leq C\left\{
 \sqrt{{\log(1/h)\over nh^2}}
 +{\log(1/h)\over n^{(1+\nu)/(2+\nu)}h^2}
 \right\},
\]
whenever \(h\downarrow0\) and \(n^{(1+\nu)/(2+\nu)}h^2/\log(1/h)\to\infty\).  The constant \(C\) depends only on the uniform class constants and \(p\).
\end{lemma}
\par\smallskip\noindent \textit{Justification.} The source-backed expected displays in SA-8.3--SA-8.4, not merely CTY Theorem~4(ii)'s probability order, justify the expected-risk argument. \textit{Gap.} The corrected locator is the current May 25, 2026 supplement: SA-8.3 (Proof of Lemma SA-3), printed pp. 32--33, and SA-8.4 (Proof of Lemma SA-4), printed pp. 33--35. \textit{Consumer.} The explicit winsorized upper proposition uses the expected Gram and score bounds.

\begin{lemma}[lem:euclidean-balls-vc]\label{lem:euclidean-balls-vc}
The collection of all closed Euclidean balls in \(\mathbb R^2\) is a VC class with VC index at most four.  Hence the fixed Euclidean metric and uniform kernel in \(\mathrm{def:cty\text{-}a1\text{-}a2\text{-}class}\) satisfy the actual VC alternative in CTY Assumption~2(iii), uniformly over every member law.
\end{lemma}
\begin{proof}
Define
\[
 \Phi(z)=(1,z_1,z_2,\lVert z\rVert_2^2)\in\mathbb R^4
\]
and, for a center \(c\in\mathbb R^2\) and radius \(r\geq0\), put
\[
 a_{c,r}=(\lVert c\rVert_2^2-r^2,-2c_1,-2c_2,1).
\]
Then
\[
 \lVert z-c\rVert_2^2-r^2=a_{c,r}^{\top}\Phi(z),
\]
so \(z\in B_2(c,r)\) if and only if \(a_{c,r}^{\top}\Phi(z)\leq0\).

Take arbitrary \(z_1,\ldots,z_4\in\mathbb R^2\) and write \(v_i=\Phi(z_i)\).  First suppose the four vectors are linearly dependent.  Choose \(\lambda_i\), not all zero, such that \(\sum_i\lambda_iv_i=0\).  The first coordinate gives \(\sum_i\lambda_i=0\), so nonzero coefficients have both signs.  Label \(z_i\) inside when \(\lambda_i>0\) and outside when \(\lambda_i<0\), assigning arbitrary labels to zero coefficients.  If a closed ball realized this labeling, then
\[
 \sum_{i=1}^4\lambda_i a_{c,r}^{\top}v_i<0,
\]
because positive-coefficient terms are nonpositive and negative coefficients multiply strictly positive outside-ball values.  This contradicts \(a_{c,r}^{\top}\sum_i\lambda_iv_i=0\).

If the \(v_i\)'s are linearly independent, there are unique \(\lambda_i\) with
\[
 \sum_{i=1}^4\lambda_iv_i=e_4,
 \qquad e_4=(0,0,0,1).
\]
Again \(\sum_i\lambda_i=0\), so both signs occur.  The same sign labeling would make the weighted sum strictly negative, but
\[
 \sum_{i=1}^4\lambda_i a_{c,r}^{\top}v_i
 =a_{c,r}^{\top}e_4=1,
\]
a contradiction.  Thus no four planar points are shattered by closed Euclidean balls.  Their VC dimension is at most three and their VC index is at most four.  Restriction to \(\mathcal X_P\) cannot increase the index.  Because Euclidean distance and \(K_\square(u)=\mathbf1\{|u|\leq1\}\) generate precisely these restricted sublevel sets, the VC clause in \(\mathrm{def:cty\text{-}a1\text{-}a2\text{-}class}\) follows uniformly.
\end{proof}
\par\smallskip\noindent \textit{Justification.} This proves, rather than assumes, the uniform-kernel VC clause required by CTY Assumption~2(iii). \textit{Gap.} The rejected extension relied only on metric equivalence and did not verify the actual VC alternative. \textit{Consumer.} The exact class definition and hard-family membership proof use this uniform finite-VC certificate.

\begin{lemma}[lem:cty-winsorization-bias]\label{lem:cty-winsorization-bias}
For \(B\geq1\), define \(\psi_B(y)=\operatorname{sign}(y)\min\{|y|,B\}\).  If \(P\in\mathcal P_{12}(p,\nu,L)\) with \(p\geq0\), \(\nu\geq2\), and \(L\geq4\), then, for both \(t\in\{0,1\}\) and every \(x\in\mathcal X_P\),
\[
 \left|\mathbb E_P\!\left[Y(t)-\psi_B\{Y(t)\}\mid X=x\right]\right|
 \leq L B^{-3}.
\]
\end{lemma}
\begin{proof}
Fix \(t\), \(x\), and \(B\geq1\).  Pointwise,
\[
 |y-\psi_B(y)|=(|y|-B)_+\leq |y|\mathbf1\{|y|>B\}.
\]
Because \(\nu\geq2\), on \(\{|y|>B\}\),
\[
 |y|\leq {|y|^{2+\nu}\over B^{1+\nu}}
 \leq {|y|^{2+\nu}\over B^3}.
\]
Conditional expectation and clause \(\mathrm{(v)}\) of \(\mathrm{def:cty\text{-}a1\text{-}a2\text{-}class}\) therefore give
\[
 \left|\mathbb E_P\!\left[Y(t)-\psi_B\{Y(t)\}\mid X=x\right]\right|
 \leq B^{-3}\mathbb E_P[|Y(t)|^{2+\nu}\mid X=x]
 \leq LB^{-3}.
\]
No smoothness of the winsorized conditional mean is used.
\end{proof}
\par\smallskip\noindent \textit{Justification.} It converts the raw fourth-or-higher conditional moment into the exact deterministic bias needed at the endpoint. \textit{Gap.} No smoothness of the winsorized conditional regression is assumed. \textit{Consumer.} The winsorized expected outer-risk proposition uses it with \(B_n=a_n^{-1/3}\).

\begin{lemma}[lem:cty-winsorized-score-maximal-bound]\label{lem:cty-winsorized-score-maximal-bound}
Let \(p\geq0\), \(\nu\geq2\), \(L\geq4\), \(P\in\mathcal P_{12}(p,\nu,L)\), \(n\geq1\), \(0<h\leq L^{-1}\), and \(B\geq1\).  Let \(\beta^*_{t,P,x}(h)\) solve the original, unwinsorized population normal equations, and define
\[
 \widehat O^B_{t,x}(h)={1\over nh^2}\sum_{i=1}^n
 \mathbf1\{D^{\pm}_{P,x,i}\in I_t\}K_\square(D^{\pm}_{P,x,i}/h)
 r_p(D^{\pm}_{P,x,i}/h)
 \left[\psi_B(Y_i)-r_p(D^{\pm}_{P,x,i}/h)^\top\beta^*_{t,P,x}(h)\right].
\]
There is a constant \(C=C(p,L)\) such that
\[
 \mathbb E_P^*\!\left[
 \sup_{t\in\{0,1\},\,x\in\mathcal B_P}
 \left\|\widehat O^B_{t,x}(h)-\mathbb E_P\widehat O^B_{t,x}(h)\right\|
 \right]
 \leq C\left\{
 \sqrt{{\log(B/h)\over nh^2}}+{B\log(B/h)\over nh^2}
 \right\}.
\]
This is a direct bounded-envelope adaptation of the expected maximal-inequality arguments in Supplemental Appendix SA-8.3 (Proof of Lemma SA-3), printed pp. 32--33, and SA-8.4 (Proof of Lemma SA-4), printed pp. 33--35, of Cattaneo, Titiunik, and Yu, not a verbatim theorem attributed to them.
\end{lemma}
\begin{proof}
Fix the displayed parameters.  On the kernel support, \(|D^{\pm}_{P,x}/h|\leq1\), so every coordinate of \(r_p(D^{\pm}_{P,x}/h)\) is bounded by one.  Clause \(\mathrm{(viii)}\) of \(\mathrm{def:cty\text{-}a1\text{-}a2\text{-}class}\) gives \(\|\Psi_{t,P,x}(h)^{-1}\|\leq L\).  The density upper bound and containment in a radius-\(h\) disk give
\[
 h^{-2}P\{|D^{\pm}_{P,x}|\leq h\}\leq \pi L.
\]
Conditional Jensen and clause \(\mathrm{(v)}\) give
\[
 \mathbb E_P[Y(t)^2\mid X]\leq L^{2/(2+\nu)}\leq L.
\]
The right side of the unwinsorized population normal equations is therefore uniformly bounded, and hence
\[
 \sup_{P,t,x,h}\|\beta^*_{t,P,x}(h)\|\leq C_0(p,L).
\]

For a coordinate \(0\leq k\leq p\), enlarge the actual summand class by allowing an independent coefficient \(b\) in the ball \(\{b\in\mathbb R^{p+1}:\|b\|\leq C_0\}\):
\[
\begin{split}
 g_{t,x,b,k}(y,z)
 ={}&\mathbf1\{D^{\pm}_{P,x}(z)\in I_t\}
 K_\square(D^{\pm}_{P,x}(z)/h)
 \left({D^{\pm}_{P,x}(z)\over h}\right)^k\\
 &\times\left[\psi_B(y)-\sum_{\ell=0}^p b_\ell
 \left({D^{\pm}_{P,x}(z)\over h}\right)^\ell\right].
\end{split}
\]
The actual class is the subclass obtained by taking \(b=\beta^*_{t,P,x}(h)\).  By \(\mathrm{lem:euclidean\text{-}balls\text{-}vc}\), the support indicators form a VC class of fixed index.  On their support the normalized radius lies in \([-1,1]\).  Closure of fixed-index VC-subgraph classes under finitely many products and sums, together with the fixed-dimensional linear parameter \(b\), gives constants \(A,v<\infty\), depending only on \(p\), such that
\[
 \sup_Q N\!\left(\epsilon\|F\|_{Q,2},\mathcal G_{t,k},L_2(Q)\right)
 \leq (A/\epsilon)^v,\qquad 0<\epsilon\leq1.
\]
Here \(F=C(p,L)B\) is a valid unscaled envelope because \(B\geq1\).  This enlarged-class argument controls the \(x\)-dependent population coefficient; it does not treat that coefficient as fixed across \(x\).

Winsorization is a contraction in absolute value, \(|\psi_B(Y)|\leq|Y|\).  The conditional second-moment bound and the fact that a summand vanishes unless \(|D^{\pm}_{P,x}|\leq h\) imply
\[
 \sup_{g\in\mathcal G_{t,k}}Pg^2\leq C(p,L)h^2.
\]
Thus the unscaled class has envelope \(U=C(p,L)B\) and standard-deviation bound \(\sigma=C(p,L)h\).  Applying the symmetrization, VC entropy-integration, and Bernstein calculation used in \(\mathrm{lem:cty\text{-}expected\text{-}local\text{-}polynomial\text{-}maximal\text{-}bounds}\) to this bounded class gives
\[
 \mathbb E_P^*\sup_{g\in\mathcal G_{t,k}}|(P_n-P)g|
 \leq C\left\{
 {h\sqrt{\log(B/h)}\over\sqrt n}
 +{B\log(B/h)\over n}
 \right\}.
\]
Indeed, the entropy ratio is \(U/\sigma\asymp B/h\).  Since \(h\leq L^{-1}\leq1/4\) and \(B\geq1\), the logarithm is positive and constants inside it are absorbed into \(C(p,L)\).  Dividing by \(h^2\), taking the maximum over the two arms, and summing the fixed \(p+1\) coordinates yields
\[
 \mathbb E_P^*\!\left[
 \sup_{t,x}\left\|\widehat O^B_{t,x}(h)-\mathbb E_P\widehat O^B_{t,x}(h)\right\|
 \right]
 \leq C(p,L)\left\{
 \sqrt{{\log(B/h)\over nh^2}}+{B\log(B/h)\over nh^2}
 \right\}.
\]
No asymptotic qualification and no joint regularity of a decision rule is used.
\end{proof}
\par\smallskip\noindent \textit{Justification.} It supplies the bounded-envelope expectation inequality with the correct entropy factor \(\log(B/h)\). \textit{Gap.} The proof enlarges the residual class to \((x,b)\) with bounded \(b\), so the \(x\)-dependent population coefficient is genuinely covered. \textit{Consumer.} The winsorized expected outer-risk proposition evaluates it at \(h_n=a_n\) and \(B_n=a_n^{-1/3}\).

\begin{lemma}[lem:cty-a1-a2-rectangle-angular-hypercube-all-orders]\label{lem:cty-a1-a2-rectangle-angular-hypercube-all-orders}
For every integer \(p\geq0\) and every \(\nu\geq2\), there is a finite \(L_0=L_0(p)\geq48\) such that, for every \(L\geq L_0\) and all sufficiently small \(\Delta>0\), the class \(\mathcal P_{12}(p,\nu,L)\) contains a hypercube \(\{P_\omega:\omega\in\{0,1\}^M\}\) with common known geometry
\[
 \mathcal X=[-3,3]^2,
 \quad \mathcal A_1=[-1,1]\times[0,2],
 \quad \mathcal A_0=\mathcal X\setminus\mathcal A_1,
 \quad \mathcal B=\partial\mathcal A_1,
\]
Euclidean distance, and \(K_\square\).  Writing \(q=p+1\geq1\), the family has \(M\geq c\Delta^{-1/q}\) disjoint disks \(S_j=B_2(x_j,w)\), with \(x_j\) on the middle of the bottom edge and \(w=A\Delta^{1/q}\), and satisfies: \(\mathrm{(a)}\) every member obeys every clause of \(\mathrm{def:cty\text{-}a1\text{-}a2\text{-}class}\); \(\mathrm{(b)}\) the cell probabilities are bit-independent and equal to \(\rho=\pi w^2/36\), the law restricted to \(S_j\) depends only on \(\omega_j\), and the law outside \(\bigcup_jS_j\) is bit-independent; \(\mathrm{(c)}\) adjacent vertices satisfy \(|\tau_{P_\omega}(x_j)-\tau_{P_{\omega^{(j)}}}(x_j)|=\Delta\); and \(\mathrm{(d)}\), conditional on \(X\in S_j\), the adjacent one-observation signed-distance laws \(Q_{j,0}\) and \(Q_{j,1}\) have the same signed-radius marginal, agree outside \(0<D^{\pm}_{P,x_j}<2C\Delta\), and obey in both orientations
\[
 \operatorname{KL}(Q_{j,0},Q_{j,1})\leq C_0{\Delta^4\over w^2}.
\]
The constants are independent of \(M\), \(\omega\), \(\Delta\), and \(n\).  When \(p=0\), the constants are chosen with \(A>2C\), so \(2C\Delta<w=A\Delta\).
\end{lemma}
\begin{proof}
Put \(q=p+1\geq1\).  Fix smooth functions \(\phi,\eta:[0,\infty)\to[0,1]\) such that \(\phi=1\) near zero, \(\phi=0\) on \([1,\infty)\), all transition endpoints are flat, \(\eta(u)=0\) for \(u\leq1\), and \(\eta(u)=1\) for \(u\geq2\).  Choose a small slope \(a>0\), then a cutoff constant \(C\) so large that \(2/(aC)\leq1/4\), and finally a smoothness constant \(A\) so large that all derivative bounds below hold and \(A>2C\).  Set
\[
 w=A\Delta^{1/q}.
\]
Place \(x_j=(s_j,0)\) on \([-1/2,1/2]\times\{0\}\) with separation at least \(4w\).  For all sufficiently small \(\Delta\), this gives \(M\asymp w^{-1}\geq c\Delta^{-1/q}\); the disks \(S_j=B_2(x_j,w)\) are disjoint, lie strictly inside \(\mathcal X=[-3,3]^2\), and avoid the corners of \(\mathcal B=\partial([-1,1]\times[0,2])\).

For \(r>0\), define
\[
 b_{\Delta,w}(r)={2\Delta\over ar}\phi(r/w)\eta(r/(C\Delta)),
 \qquad b_{\Delta,w}(0)=0.
\]
For \(x\in S_j\), write \(x-x_j=r_j(\cos\alpha_j,\sin\alpha_j)\), interpret every summand as zero off its disk, and put
\[
 f_\omega(x)={1\over36}\left[1+\sum_{j=1}^M\omega_jb_{\Delta,w}(r_j)\cos\alpha_j\right],
\]
\[
 \mu_{0,\omega}(x)={1\over2},
 \qquad
 \mu_{1,\omega}(x)={1\over2}+ax_1-\Delta\sum_{j=1}^M\omega_j\phi(r_j/w).
\]
Conditional on \(X=x\), take \(Y(0)\) and \(Y(1)\) conditionally independent Bernoulli variables with these means, and impose the fixed deterministic assignment and consistency equations.

The density perturbation vanishes for \(r\leq C\Delta\) and smoothly before \(r=w\).  The choice \(A>2C\) gives \(2C\Delta<w\) when \(q=1\); for \(q>1\) the same inequality holds for all sufficiently small \(\Delta\).  Wherever the perturbation is nonzero,
\[
 |b_{\Delta,w}(r)|\leq {2\over aC}\leq{1\over4}.
\]
Disjointness gives \(1/48\leq f_\omega\leq5/144\).  Moreover,
\[
 \int_0^w b_{\Delta,w}(r)r\,dr\int_0^{2\pi}\cos\alpha\,d\alpha=0
\]
on every disk, so \(\int_{\mathcal X}f_\omega=1\).  Thus the density is continuous, normalized, positive on the entire square, and satisfies the envelope whenever \(L\geq48\).

For every derivative order \(0\leq k\leq q\), scaling and disjointness yield
\[
 \left\|D^k\left[\Delta\sum_j\omega_j
 \phi(\|\,\cdot-x_j\|_2/w)\right]\right\|_\infty
 \leq C_k\Delta w^{-k}
 =C_kA^{-k}\Delta^{1-k/q}.
\]
At \(k=q\) this is \(C_qA^{-q}\), and at smaller orders it vanishes with \(\Delta\).  This proves the derivative and Lipschitz bounds through order \(p\).  At \(p=0,q=1\), the gradient is bounded by \(C_1/A\), so the regression itself has the required Lipschitz constant.  Taking \(a\) and then \(\Delta\) small places both Bernoulli probabilities in \([1/4,3/4]\).  Their variances lie in \([3/16,1/4]\), and their raw conditional \((2+\nu)\)-moments are at most one.  Smooth extension holds because the disks are strictly interior and the affine part is global.

The assignment rectangle has perimeter eight and distance at least one from the support boundary.  If \(0<h\leq L^{-1}\) and \(L\) is sufficiently large, then at every edge or corner each arm contains a sector of radius \(h\) and opening angle at least \(\pi/2\).  Hence
\[
 {1\over h^2}\int_{\mathcal A_t}
 K_\square(\{(2t-1)\|z-x\|_2\}/h)\,dz\geq{\pi\over8}.
\]
For \(0<s\leq L^{-1}\), the radius-\(s\) arc in either arm has length between \((\pi/2)s\) and \(2\pi s\); multiplication by \(1/48\leq f_\omega\leq5/144\) proves that the radial Hausdorff integrals are finite and strictly positive.  For \(v\in\mathbb R^{p+1}\), polar coordinates in such a sector and the density lower bound give
\[
 v^\top\Psi_{t,P_\omega,x}(h)v
 \geq c\int_0^1\{v^\top r_p((2t-1)u)\}^2u\,du.
\]
The integral is positive for \(v\ne0\), including \(p=0\).  Compactness of the Euclidean unit sphere in \(\mathbb R^{p+1}\) gives a positive uniform minimum.  Enlarging \(L_0(p)\) makes this minimum at least \(L^{-1}\) and also verifies the support, density, variance, boundary-length, and local-mass envelopes.  The Euclidean-ball VC clause follows from \(\mathrm{lem:euclidean\text{-}balls\text{-}vc}\).  Thus every clause of \(\mathrm{def:cty\text{-}a1\text{-}a2\text{-}class}\) holds.

Angular integration gives
\[
 P_\omega(X\in S_j)={|S_j|\over36}={\pi w^2\over36}=\rho.
\]
Disjoint support gives bit locality.  At \(x_j\),
\[
 \tau_{P_\omega}(x_j)=as_j-\Delta\omega_j,
\]
so adjacent targets differ by \(\Delta\).

Inside \(S_j\), positive signed radius corresponds to the upper semicircle.  Using
\[
 \int_0^\pi\cos\alpha\,d\alpha=0,
 \qquad \int_0^\pi\cos^2\alpha\,d\alpha={\pi\over2},
\]
shows first that the signed-radius marginal is bit-independent and then that angular averaging on the treatment half gives
\[
 q_{1,j}(r)-q_{0,j}(r)
 =-\Delta\phi(r/w)\{1-\eta(r/(C\Delta))\}.
\]
On the control half the observed Bernoulli mean is \(1/2\), so its radially compressed outcome law is bit-independent as well.  Hence the two conditional laws agree outside \(0<r<2C\Delta\), an interval contained in \((0,w)\).  Since all Bernoulli probabilities lie in \([1/4,3/4]\), in either KL orientation,
\[
 \operatorname{kl}(u,v)\leq{(u-v)^2\over v(1-v)}
 \leq C_1\Delta^2.
\]
Conditional on the disk, each signed half-radius has density \(r/w^2\).  Disintegration therefore yields
\[
 \operatorname{KL}(Q_{j,0},Q_{j,1})
 \leq C_1\Delta^2\int_0^{2C\Delta}{r\over w^2}\,dr
 \leq C_0{\Delta^4\over w^2}.
\]
All constants are independent of \(M\), \(\omega\), \(\Delta\), and \(n\), proving the lemma for every \(p\geq0\).
\end{proof}
\par\smallskip\noindent \textit{Justification.} This is the same-law-class least-favorable family for every polynomial order, including local constants. \textit{Gap.} The endpoint uses \(A>2C\) to keep the cancellation annulus inside each disk and retains the full first-derivative bound. \textit{Consumer.} The all-order point-indexed converse uses its local divergence, target separation, equal masses, and locality.

\begin{theorem}[thm:cty-a1-a2-point-indexed-log-converse-all-orders]\label{thm:cty-a1-a2-point-indexed-log-converse-all-orders}
For every integer \(p\geq0\) and every \(\nu\geq2\), there is a finite \(L_0=L_0(p)\) such that, for every \(L\geq L_0\), there is \(c=c(p,\nu,L)>0\) satisfying
\[
 \liminf_{n\to\infty}\left({n\over\log n}\right)^{1/4}R^{12,\pm}_n(p,\nu,L)\geq c.
\]
The conclusion holds for arbitrary law-independent known-geometry point-indexed families with Borel fixed sections, under outer expectation and without joint regularity in the boundary index.  The hard subclass has one fixed support, one fixed known assignment rectangle, and one fixed known interior treatment boundary.
\end{theorem}
\begin{proof}
Fix \(p\geq0\), \(\nu\geq2\), and \(L\geq L_0(p)\), and put \(q=p+1\geq1\).  Let
\[
 \Delta_n=\gamma a_n,
 \qquad w_n=A\Delta_n^{1/q},
\]
where \(\gamma>0\) will be chosen small.  Apply \(\mathrm{lem:cty\text{-}a1\text{-}a2\text{-}rectangle\text{-}angular\text{-}hypercube\text{-}all\text{-}orders}\).  It gives \(M_n\asymp w_n^{-1}\to\infty\), common cell probability \(\rho_n=\pi w_n^2/36\), target separation \(\Delta_n\), and adjacent conditional cell divergence at most \(C_0\Delta_n^4/w_n^2\).

Put independent uniform bits on the vertices.  For the \(n\) observations, let \(J_i\) record the cell containing \(X_i\), with label zero outside all cells.  Equal cell masses and locality make the label vector bit-independent.  Conditional on that vector, the cell blocks are independent, block \(j\) depends only on bit \(j\), and block zero is bit-independent.  Marginally \(N_j=\sum_i\mathbf1\{J_i=j\}\sim\operatorname{Bin}(n,\rho_n)\).  Exponential Markov inequality at \(s=\log2\), followed by \(1+u\leq e^u\), gives
\[
\begin{split}
 \Pr(N_j>2n\rho_n)
 &\leq e^{-2sn\rho_n}(1-\rho_n+\rho_ne^s)^n\\
 &\leq\exp\{-n\rho_n(2\log2-1)\}
 \leq\exp(-n\rho_n/3).
\end{split}
\]
A union bound therefore gives
\[
 \Pr\!\left(\max_jN_j>2n\rho_n\right)
 \leq M_n\exp(-n\rho_n/3).
\]
Now
\[
 n\rho_n\asymp n\Delta_n^{2/q}
 =n^{1-1/(2q)}(\log n)^{1/(2q)},
 \qquad \log M_n=O(\log n).
\]
Their ratio diverges for every \(q\geq1\); at \(q=1\) it is of order \(n^{1/2}(\log n)^{-1/2}\).  Consequently the good-count event
\[
 G_n=\{\max_jN_j\leq2n\rho_n\}
\]
has probability tending to one.

Fix a label vector in \(G_n\).  For decoder \(j\), reveal the complete off-cell raw blocks and retain from its own block only the ordered outcome--signed-distance observations.  Together with the fixed label vector, these data reconstruct the compressed sample at \(x_j\), so the fixed Borel section and its midpoint decoder are measurable with respect to \((S_j,Z_{-j},Z_0)\).  Conditional tensorization gives, in both orientations,
\[
 \operatorname{KL}\{\mathcal L(S_j\mid\Omega_j=0,J_{1:n}),
 \mathcal L(S_j\mid\Omega_j=1,J_{1:n})\}
 \leq N_jC_0{\Delta_n^4\over w_n^2}
 \leq C_1n\Delta_n^4.
\]
Since \(n\Delta_n^4=\gamma^4\log n\) and
\[
 \log M_n={1\over4q}\log n+O(\log\log n),
\]
choose \(\gamma>0\) so small that the divergence is at most \(\kappa\log M_n\) for a fixed \(\kappa<1\).  Conditional on labels in \(G_n\), \(\mathrm{lem:coordinatewise\text{-}overlap\text{-}direct\text{-}product}\) yields
\[
 \Pr\{\text{some coordinate decoder errs}\mid J_{1:n}\}
 \geq {1\over2}\left[1-\exp\{-M_n^{1-\kappa}/2\}\right].
\]
Averaging over the bit-independent labels and the uniform vertex prior gives probability \(1/2-o(1)\) of at least one error.

At \(x_j\) the two possible target values are separated by \(\Delta_n\).  An erroneous midpoint decoder therefore implies that the measurable finite-coordinate loss is at least \(\Delta_n/2\), whence
\[
 2^{-M_n}\sum_\omega
 \mathbb E_{P_\omega}\max_{j\leq M_n}
 |\widehat\tau_{n,P_\omega}(x_j)-\tau_{P_\omega}(x_j)|
 \geq {\Delta_n\over2}\{1/2-o(1)\}.
\]
The finite maximum is measurable because every fixed-\((G,x_j)\) section is Borel.  It is pointwise bounded by the possibly nonmeasurable full-boundary supremum.  Monotonicity in the definition of outer expectation therefore gives the same lower bound for the outer expectation of that supremum.  Every vertex belongs to \(\mathcal P_{12}(p,\nu,L)\).  Taking the supremum over \(P\), the infimum over the arbitrary point-indexed rule, dividing by \(a_n\), and taking the lower limit proves the assertion with a positive constant proportional to \(\gamma\).  The construction has the fixed common geometry displayed in the hypercube lemma, and no joint regularity in the boundary index enters the proof.
\end{proof}
\par\smallskip\noindent \textit{Justification.} It proves the lower half of the second co-main \(\mathcal P_{12}\) frontier for every \(p\geq0\), covering arbitrary law-independent known-geometry point-indexed families with Borel fixed sections on one fixed support, assignment rectangle, and interior boundary. \textit{Gap.} This supplies the outer-expected minimax converse absent from CTY's identification, bias, and probability-order results; at \(q=1\), the expected cell count still dominates the logarithm of the packing size. \textit{Consumer.} It is the lower half of the second co-main matched \(\mathcal P_{12}\) theorem.

\begin{proposition}[prop:cty-a1-a2-winsorized-expected-outer-upper]\label{prop:cty-a1-a2-winsorized-expected-outer-upper}
For every integer \(p\geq0\), every \(\nu\geq2\), and every \(L\geq4\), the explicit winsorized, Gram-stabilized estimator in \(\mathrm{def:cty\text{-}stabilized\text{-}local\text{-}polynomial\text{-}estimator}\), with \(h_n=a_n=(\log n/n)^{1/4}\) and \(B_n=a_n^{-1/3}\), satisfies, for all sufficiently large \(n\),
\[
 \sup_{P\in\mathcal P_{12}(p,\nu,L)}
 \mathbb E_P^*\!\left[
 \sup_{x\in\mathcal B_P}
 |\widehat\tau^{12,\mathrm{LP}}_{n,h_n}(x)-\tau_P(x)|
 \right]
 \leq Ca_n.
\]
The expectation theorem is derived by comparing the winsorized empirical score directly with the original population coefficient and by a bounded-envelope adaptation of the CTY supplement's expected maximal-inequality proof; it is not attributed verbatim to CTY.
\end{proposition}
\begin{proof}
Fix an integer \(p\geq0\), numbers \(\nu\geq2\) and \(L\geq4\), and a law \(P\in\mathcal P_{12}(p,\nu,L)\).  Temporarily write \(h=h_n\) and \(B=B_n\).  For \(t\in\{0,1\}\) and \(x\in\mathcal B_P\), let \(\beta^*_{t,P,x}(h)\) solve the original, unwinsorized population normal equations, and define
\[
 D_h=\sup_{t\in\{0,1\},\,x\in\mathcal B_P}
 \lVert\widehat\Psi_{t,x}(h)-\Psi_{t,P,x}(h)\rVert,
 \qquad
 \mathcal G_h=\{D_h\leq(2L)^{-1}\}.
\]
All expectations of possibly nonmeasurable suprema below are outer expectations.

On \(\mathcal G_h\), clause \(\mathrm{(viii)}\) of \(\mathrm{def:cty\text{-}a1\text{-}a2\text{-}class}\) and Weyl's eigenvalue inequality give, uniformly in \((t,x)\),
\[
 \lambda_{\min}\{\widehat\Psi_{t,x}(h)\}
 \geq \lambda_{\min}\{\Psi_{t,P,x}(h)\}-D_h
 \geq L^{-1}-(2L)^{-1}=(2L)^{-1}.
\]
Thus the stabilization branch does not activate and \(\lVert\widehat\Psi_{t,x}(h)^{-1}\rVert\leq2L\).  The empirical normal equations and the definition of \(\widehat O^B_{t,x}(h)\) imply
\[
 \widehat\beta^B_{t,x}(h)-\beta^*_{t,P,x}(h)
 =\widehat\Psi_{t,x}(h)^{-1}\widehat O^B_{t,x}(h),
\]
and hence
\[
 \sup_{t,x}\lVert\widehat\beta^B_{t,x}(h)-\beta^*_{t,P,x}(h)\rVert
 \leq 2L\sup_{t,x}\lVert\widehat O^B_{t,x}(h)\rVert
 \quad\text{on }\mathcal G_h.
\]

Subtracting the unwinsorized population normal equations from the definition of the winsorized score yields
\[
 \mathbb E_P\widehat O^B_{t,x}(h)
 ={1\over h^2}\mathbb E_P\!\left[
 \mathbf1\{D^{\pm}_{P,x}\in I_t\}K_\square(D^{\pm}_{P,x}/h)
 r_p(D^{\pm}_{P,x}/h)\{\psi_B(Y)-Y\}
 \right].
\]
On the observed arm, clause \(\mathrm{(i)}\) gives \(Y=Y(t)\).  Each coordinate of the polynomial basis is bounded by one on the kernel support.  The conditional bound in \(\mathrm{lem:cty\text{-}winsorization\text{-}bias}\), the density bound \(f_P\leq L\), and containment of the kernel support in the radius-\(h\) disk give
\[
\begin{split}
 \sup_{t,x}\lVert\mathbb E_P\widehat O^B_{t,x}(h)\rVert
 &\leq {CB^{-3}\over h^2}
 \sup_{x\in\mathcal B_P}P\{\lVert X-x\rVert_2\leq h\}\\
 &\leq {CB^{-3}\over h^2}\,L\pi h^2
 \leq CB^{-3}.
\end{split}
\]
The centered bound in \(\mathrm{lem:cty\text{-}winsorized\text{-}score\text{-}maximal\text{-}bound}\) and the triangle inequality for outer expectation now show
\[
 \mathbb E_P^*\sup_{t,x}\lVert\widehat O^B_{t,x}(h)\rVert
 \leq C\left\{B^{-3}+\sqrt{{\log(B/h)\over nh^2}}
 +{B\log(B/h)\over nh^2}\right\}.
\]

At the bandwidth used in this proposition,
\[
 h=h_n=a_n=\left({\log n\over n}\right)^{1/4}\downarrow0,
 \qquad
 nh_n^2=\sqrt{n\log n}\longrightarrow\infty.
\]
Consequently, the sequence-qualified conclusion of \(\mathrm{lem:cty\text{-}uniform\text{-}first\text{-}order\text{-}bias}\) applies for all sufficiently large \(n\) and gives
\[
 \sup_{x\in\mathcal B_P}\left|
 e_1^\top\{\beta^*_{1,P,x}(h_n)-\beta^*_{0,P,x}(h_n)\}-\tau_P(x)
 \right|\leq Ch_n.
\]
The identification of the causal target with the signed-distance limits is the equality in \(\mathrm{lem:cty\text{-}distance\text{-}identification}\); it is distinct from the population regression coefficient used as the approximation device.  Clause \(\mathrm{(iii)}\) gives \(|\mu_{t,P}(x)|\leq L\), and therefore \(|\tau_P(x)|\leq2L\).  Projection onto \([-2L,2L]\) is one-Lipschitz and fixes \(\tau_P(x)\).  On \(\mathcal G_h^c\), both the projected estimator and the target belong to this interval, so their absolute difference is at most \(4L\).  Combining these observations pointwise gives
\[
 \sup_{x\in\mathcal B_P}
 |\widehat\tau^{12,\mathrm{LP}}_{n,h}(x)-\tau_P(x)|
 \leq Ch+C\sup_{t,x}\lVert\widehat O^B_{t,x}(h)\rVert
 +4L\mathbf1\{D_h>(2L)^{-1}\}.
\]
Since \(\mathbf1\{D_h>(2L)^{-1}\}\leq2LD_h\), the Gram bound in \(\mathrm{lem:cty\text{-}expected\text{-}local\text{-}polynomial\text{-}maximal\text{-}bounds}\) implies
\[
 \mathbb E_P^*\mathbf1\{D_h>(2L)^{-1}\}
 \leq C\sqrt{{\log(1/h)\over nh^2}}
 \leq C\sqrt{{\log(B/h)\over nh^2}},
\]
where the last inequality uses \(B\geq1\).  Monotonicity and subadditivity of outer expectation consequently yield, uniformly over \(P\in\mathcal P_{12}(p,\nu,L)\),
\[
 \mathbb E_P^*\!\left[
 \sup_{x\in\mathcal B_P}
 |\widehat\tau^{12,\mathrm{LP}}_{n,h}(x)-\tau_P(x)|
 \right]
 \leq C\left\{h+B^{-3}+\sqrt{{\log(B/h)\over nh^2}}
 +{B\log(B/h)\over nh^2}\right\}.
\]

Continue with \(h=h_n=a_n\) and set \(B=B_n=a_n^{-1/3}\).  Separately from the bias-interface limits verified above, the regime required by the established maximal inequalities holds because
\[
 {n^{(1+\nu)/(2+\nu)}h_n^2\over\log(1/h_n)}
 \asymp {n^{\nu/\{2(2+\nu)\}}\over\sqrt{\log n}}
 \longrightarrow\infty.
\]
Moreover,
\[
 \log(B_n/h_n)={4\over3}\log(1/h_n)
 ={1\over3}(\log n-\log\log n),
\]
so
\[
 h_n=a_n,\qquad B_n^{-3}=a_n,
 \qquad \sqrt{{\log(B_n/h_n)\over nh_n^2}}=O(a_n),
\]
while
\[
 {B_n\log(B_n/h_n)\over nh_n^2}
 =O\!\left(n^{-5/12}(\log n)^{5/12}\right)=o(a_n).
\]
All constants depend only on \((p,\nu,L)\) and are uniform in \(P\) and \(n\).  Taking the supremum over \(P\in\mathcal P_{12}(p,\nu,L)\) proves the asserted bound for all sufficiently large \(n\).  Finally, \(\mathrm{def:cty\text{-}stabilized\text{-}local\text{-}polynomial\text{-}estimator}\) supplies this rule as one explicit, law-independent member of \(\mathcal T_n^{12,\pm}\).
\end{proof}
\par\smallskip\noindent \textit{Justification.} It proves the upper half of the second co-main \(\mathcal P_{12}\) frontier with one explicit, law-independent winsorized, clipped, Gram-stabilized estimator under outer-expected uniform loss for every \(L\geq4\). \textit{Gap.} CTY Theorem~4(ii) gives probability-order stochastic control, not this common-law-class expected-risk achievability statement.  Direct comparison with the original population coefficient avoids imposing unproved smoothness after winsorization. \textit{Consumer.} It is the upper half of the second co-main matched theorem and remains valid throughout the broader \(L\geq4\) range even when the matching lower construction is not asserted.

\begin{theorem}[thm:cty-a1-a2-winsorized-matched-frontier]\label{thm:cty-a1-a2-winsorized-matched-frontier}
For every integer \(p\geq0\), every \(\nu\geq2\), and every \(L\geq L_0(p)\), there are constants \(0<c\leq C<\infty\) such that
\[
 0<c\leq\liminf_{n\to\infty}{R^{12,\pm}_n(p,\nu,L)\over a_n}
 \leq\limsup_{n\to\infty}{R^{12,\pm}_n(p,\nu,L)\over a_n}\leq C<\infty.
\]
Thus the outer-expected minimax rate over the same exact displayed \(L\)-uniformized Euclidean-distance, uniform-kernel CTY Assumptions~1--2 plus Theorem-2-envelope class is \(a_n=(\log n/n)^{1/4}\) throughout \(q=p+1\geq1\) and \(\nu\geq2\).  The upper bound is attained by the explicit winsorized, Gram-stabilized degree-\(p\) signed-distance local-polynomial estimator with \(h_n=a_n\) and \(B_n=a_n^{-1/3}\); the lower bound holds over all law-independent known-geometry point-indexed Borel sections.
\end{theorem}
\begin{proof}
Fix an integer \(p\geq0\) and a number \(\nu\geq2\).  Enlarge the finite threshold in \(\mathrm{thm:cty\text{-}a1\text{-}a2\text{-}point\text{-}indexed\text{-}log\text{-}converse\text{-}all\text{-}orders}\), if necessary, by replacing it with its maximum with \(4\); this preserves that theorem.  Fix \(L\geq L_0(p)\).  The point-indexed converse gives a constant \(c_0=c_0(p,\nu,L)>0\) such that
\[
 \liminf_{n\to\infty}{R_n^{12,\pm}(p,\nu,L)\over a_n}\geq c_0.
\]
This lower bound uses the same law class, the same outer-expectation risk, and the full law-independent known-geometry point-indexed decision class with Borel fixed sections.

Because \(L\geq4\), the explicit winsorized, Gram-stabilized estimator is covered by \(\mathrm{prop:cty\text{-}a1\text{-}a2\text{-}winsorized\text{-}expected\text{-}outer\text{-}upper}\).  It is a single law-independent member of \(\mathcal T_n^{12,\pm}\).  Evaluating the infimum in \(\mathrm{def:cty\text{-}a1\text{-}a2\text{-}outer\text{-}risk}\) at this estimator gives, for all sufficiently large \(n\),
\[
\begin{split}
 R_n^{12,\pm}(p,\nu,L)
 &\leq \sup_{P\in\mathcal P_{12}(p,\nu,L)}
 \mathbb E_P^*\!\left[
 \sup_{x\in\mathcal B_P}
 |\widehat\tau^{12,\mathrm{LP}}_{n,a_n}(x)-\tau_P(x)|
 \right]\\
 &\leq C_0a_n
\end{split}
\]
for a finite \(C_0=C_0(p,\nu,L)\).  Since \(a_n>0\), division by \(a_n\) and passage to the upper limit yield
\[
 \limsup_{n\to\infty}{R_n^{12,\pm}(p,\nu,L)\over a_n}\leq C_0<\infty.
\]
For every real sequence, its lower limit does not exceed its upper limit.  Taking \(c=c_0/2\) and \(C=\max\{C_0,c\}\) proves
\[
 0<c\leq\liminf_{n\to\infty}{R_n^{12,\pm}(p,\nu,L)\over a_n}
 \leq\limsup_{n\to\infty}{R_n^{12,\pm}(p,\nu,L)\over a_n}
 \leq C<\infty.
\]
Finally, \(a_n=(\log n/n)^{1/4}\) by \(\mathrm{def:frontier\text{-}rates}\), and \(q=p+1\geq1\).  The upper witness is valid already for every \(L\geq4\); the matched conclusion uses \(L\geq L_0(p)\) precisely because that is the threshold required by the converse.
\end{proof}
\par\smallskip\noindent \textit{Justification.} This is the second unrelegated co-main theorem.  It is important because it turns CTY's identification, sequence-scoped bias, and probability-order stochastic ingredients into an outer-expected minimax equality on one exact displayed \(\mathcal P_{12}(p,\nu,L)\) law, risk, and decision triple. \textit{Gap.} CTY Theorem~1 identifies the curve, Theorem~2 supplies the sequence-scoped first-order bias input, Theorem~4(ii) supplies probability-order stochastic control, and Theorem~6 concerns the distinct \(\mathcal P_{\mathrm{NP}}\) support-boundary problem.  CTY do not state this common-class outer-expected lower/upper equality; the fixed-known-interior-boundary converse over arbitrary point-indexed Borel sections and the explicit winsorized, Gram-stabilized expected-risk witness close that mismatch. \textit{Consumer.} It is the paper's second co-main result, complementary to and non-substitutable for the first co-main full-\(\mathcal P_{\mathrm{NP}}\) converse.

\section{Interpretation}

The two results isolate different frontiers and neither substitutes for the other.  On full \(\mathcal P_{\mathrm{NP}}\), the logarithm is an information-loss phenomenon for simultaneous support-boundary recovery from unsigned-distance samples, not an artifact of imposing one common estimator map; the conclusion is converse-only.  On the narrower and structurally different \(\mathcal P_{12}\) causal model, signed-distance identification and uniform Gram geometry permit an explicit estimator, while a fixed-geometry hypercube shows that no arbitrary point-indexed Borel family improves the order.  Their combination yields the separate matched outer-expected frontier only on \(\mathcal P_{12}\).

\section*{Technical internal limitation (diagnostic only)}

The full-\(\mathcal P_{\mathrm{NP}}(L,q)\) theorem is only a lower bound and does not inherit the \(\mathcal P_{12}\) upper result.  The matched theorem is confined to the exact displayed \(L\)-uniformized Euclidean-distance, uniform-kernel CTY Assumptions~1--2 plus Theorem-2-envelope class; no broader metric or kernel regime and no exact asymptotic constant are claimed.  Its expected-risk upper holds for \(L\geq4\), whereas its available hard subclass requires \(L\geq L_0(p)\).  Outer expectation is retained because arbitrary point-indexed families need not be jointly measurable in the boundary index.

\section{Honest scope}

The first co-main contribution is the \((\log n/n)^{1/4}\) converse on the exact full \(\mathcal P_{\mathrm{NP}}(L,q)\) class for integer \(q\geq1\) and \(L\geq4\), covering arbitrary law-independent point-indexed Borel sections and therefore the literal CTY common-map class; it claims no full-class achievability.  The second co-main contribution is the explicit winsorized expected-risk upper on the distinct exact displayed \(\mathcal P_{12}(p,\nu,L)\) specialization for every \(L\geq4\), together with the matched outer-expected minimax equality for \(L\geq L_0(p)\).  The note does not claim broader CTY metric/kernel coverage, transfer of this upper bound to full \(\mathcal P_{\mathrm{NP}}\), or an exact minimax constant.

\begin{thebibliography}{99}

  \bibitem{CattaneoTitiunikYu2026Distance} Cattaneo, Matias D., Rocio Titiunik, and Ruiqi Rae Yu. 2026. Estimation and Inference in Boundary Discontinuity Designs: Distance-Based Methods. Journal of Econometrics 256, article 106266. doi:10.1016/j.jeconom.2026.106266. Supplemental Appendix, May 25, 2026, \url{https://mdcattaneo.github.io/papers/Cattaneo-Titiunik-Yu\_2026\_JOE--Supplement.pdf}.

\end{thebibliography}

\end{document}